\chapter{知识的更新和维护}
\label{chap:ket-update-maintenance}

一个仍需识别原有物种的花卉系统，在春季照片上错误较少，进入雨季后却频繁误判。我们需要查明：是原来少见的阴影照片变多了，还是同样的记录特征已不足以判断物种？系统附带的植物资料问答也可能需要改动，例如研究机构更换了负责人，当前答案应更新，任命前的历史答案却应保留。若一条记录被依法撤回，还要处理它已经造成的训练影响，不能只让系统停止输出某个名字。这些请求都针对已有系统，但要求改变的对象和验收依据不同。

第~\ref{chap:ket-reuse-selection}章解决了是否复用以及选择什么知识的问题。本章接着考察已经承担服务责任的系统：新证据到来后，怎样吸收需要改变的内容，又不破坏仍有效的能力？为此，更新目标先规定哪些响应应变、哪些应留，证据准备确定我们凭什么作出这种判断，操作把证据转成系统变化，验证再决定是否接受候选版本。多轮更新还需要保存请求、预算和恢复状态。若历史回答被误改，就返回范围判断；若标签不能辨认失配，就补证据。面向未见环境的预防性训练也属于这条路线，但训练和选模不能提前使用未来目标数据。

以下沿用前文的联合分布$\mathcal P_S,\mathcal P_T$。风险$R_S,R_T$表示相应分布下的平均任务损失，$\hat R$表示样本上的平均损失；比较两域风险时，输入输出语义与损失须一致。模型记为$f_{\bm{\theta}}$，输入为$\bm{x}$，表示为$\bm{z}$，上标$^{(t)}$标记学习阶段或迭代，下标$k$标记任务$\mathcal T_k$。理论卷\BookChapterRef{chap:transfer-learning-theory}{卷3《基础、理论与可信性》第9章“迁移学习的可学习性”}已经说明：没有覆盖或连接两域标签规律的依据，无标签输入不能确定任意新规律。本章据此讨论能实施的更新，不把可计算的代理目标当成可学习性保证。输出语义或成功标准扩展时，第~\ref{sec:ket-transfer-output-semantics}节先确定新任务；本章承担扩展后的保留和维护。现成模型如何整合、多个任务如何共同训练，分别由第~\ref{chap:ket-integration}章和第~\ref{chap:ket-sharing}章展开。

在选择算法之前，可以把一次维护请求放入表~\ref{tab:ket-update-route-matrix}的四项检查中：要达到什么目标，哪个载体需要变化，允许访问什么信息，以及证据按怎样的时间协议到来。例如，雨季适应若只有当前无标签流，就不能调用需要旧照片和即时标签的操作；负责人修正若带有效时间，就不能无条件覆盖所有历史回答。表中各列分别限制操作的一个方面，各行描述主要维护责任。一项请求可能同时涉及几行，并不需要被归入唯一类别。

\begin{table*}[htbp]
\small
\begin{tabularx}{\linewidth}{@{}p{23mm}XXXX@{}}
\toprule
更新目标 & 优先检查的载体 & 关键权限 & 时间协议 & 路线选择问题 \\
\midrule
环境适应 & 预处理、表示、参数与统计状态 & 当前输入；标签或源接口按协议声明 & 离线批量、测试时或严格在线 & 变化是否真的破坏任务规律，代理信号能否支持更新 \\
事实修正 & 外部记录、索引、缓存或参数行为 & 有效请求、来源、依赖与历史查询 & 生效区间和连续版本 & 改哪里能同时满足新事实、历史语义与局部性 \\
持续吸收 & 样本记忆、参数、模块和优化状态 & 历史访问、容量与反馈时机 & 阶段序列、环境回归或无边界流 & 怎样兼顾当前吸收、旧能力与后续可塑性 \\
训练影响消除 & 原记录及派生物、参数与恢复快照 & 撤回集合、剩余数据和规定学习过程 & 请求到达后持续有效 & 外部删除与参数影响消除分别需要什么证据 \\
\bottomrule
\end{tabularx}
\caption{选择更新路线时需要共同检查的目标、载体、权限与时间协议。各行按主要维护责任组织，不构成互斥分类。在线事实修正可以同时改变记录、缓存与参数；每项操作仍须满足该请求的授权和有效时间条件。}
\label{tab:ket-update-route-matrix}
\end{table*}

\section{知识更新目标的确定}
\label{sec:ket-update-goals}

雨季错误增多时，先训练一个新版本还不能说明改对了什么。我们需要先检查错误来自数据、观测接口还是预测规律，再划定允许改变的内容，最后约定仍需服务的能力和可用资源。诊断决定是否更新，范围决定哪里可以变，保留要求决定候选版本能否接受。这三项判断相互制约，却不能由一次训练损失下降同时回答。

\subsection{更新需求的辨认}
\label{sec:ket-update-needs}

先保持照片的记录方式、物种集合和错误计分规则不变，只考察照片及标签怎样共同出现。若两域在共同覆盖区域内满足$p_T(y\mid\bm{x})=p_S(y\mid\bm{x})$，而输入分布不同，就得到\BookSectionRef{sec:distribution-shift}{卷3《基础、理论与可信性》第9章“迁移学习的可学习性”第9.1.2节“分布偏移与可识别性”}的\term{协变量偏移}（covariate shift）：同一照片对应的标签概率不变，变化的是看到哪些照片的频率。若保持$p_T(\bm{x}\mid y)=p_S(\bm{x}\mid y)$而改变$p_T(y)$，则属于\term{标签偏移}（label shift）：每个物种内部的照片分布保持，变化的是物种比例。后者通常会改变后验$p(y\mid\bm{x})$，因此两种假设不能互换。

观察到$p_T(y\mid\bm{x})$改变后，还要检查类条件分布是否保持。两种近似物种若在已保存的压缩特征上难以区分，只改变物种比例便可能改变后验；这种情况仍可符合标签偏移。若相机处理又丢掉了细微纹理，物种内部的特征分布也会变化，只校正比例就没有足够依据。这里的条件分布始终相对于\emph{实际记录的输入}定义，物种的生物学定义不变并不保证它不变。若标注规则重新划分了物种，则须先核对输出语义，而不能把所有差异都解释为固定任务下的统计偏移。

如果彩色相机换成深度相机，我们改变的不只是样本频率，还改变了取得信息的接口。这种\term{观测空间变化}（change of observation space）可以与统计偏移同时发生。RGB记录红、绿、蓝通道强度，深度图记录距离；即便整理成同长向量，对应坐标也没有相同含义。相反，同一照片改用另一种无损编码，文件长度变了，可判别信息未必变。本章以RGB到深度图的教学情形补充概述中的传感器更换，考察植株和器官怎样对应，以及新观测是否还包含区分物种所需的花色；它不与多光谱测量混作同一输入。

把三种部署变化放在一起，区别便落到了可检查的条件上。温室转露天仍记录照片，却可能同时改变背景、遮挡和物种比例，地点名称不能证明协变量偏移假设成立。生长阶段若使同样的记录特征对应不同物种概率，变化发生在条件预测关系。深度相机则要求重新建立观测与物种之间的联系，不能直接沿用RGB坐标。仿真到现实若只改变成像条件，也可逐项作这些检查；若还改变动力学、奖励或交互后果，就需要\BookVolumeRef{vol:reinforcement-learning}{卷7《学习范式：强化学习》}中的专门建模，不能由图像适应结论代替。

变化告警本身不能决定训练。设识别器只读取花瓣形态，背景由灰变绿，但形态、标签关系和观测质量都保持，原规则仍可能适用，此时不更新应留在候选中。反过来，准确率下降也可能来自错标、距离单位改变、解码故障或评价权重变化。先核对这些环节，才能判断是否需要前文的\term{领域自适应}（亦称领域适应，domain adaptation）：利用当前域的许可证据调整系统，以改善该域上的原任务预测。它是待检验的解决路线，不是告警之后必然成功的操作。

\subsection{修改内容与适用范围的确定}
\label{sec:ket-update-scope}

照片适应要从一批观测中辨认失配，负责人变更却可以由一份任命请求触发。后者必须把“改成谁”与“何时、对哪个机构生效”一起保存，才能生成训练和验证条件。将一条定向请求记为
\begin{equation}
e=(s,r,o_{\mathrm{old}},o_{\mathrm{new}},c,[\tau_0,\tau_1),a),
\label{eq:ket-update-request}
\end{equation}
其中$s$标明主体，$r$标明关系，$o_{\mathrm{old}},o_{\mathrm{new}}$分别是修改前与期望记录，$c$补充适用条件，$a$保存依据。有效区间$[\tau_0,\tau_1)$包含起点而不含终点，$\tau$指日历时间，与迭代编号$t$不同。请求只有在主体、关系、时间及其他条件匹配时才适用。这些字段共同确定应改变的响应，单靠关键词替换无法保留这层含义。

以下机构、人物和日期均为教学构造。植物研究机构甲所以独立标识区别于同名机构，人员也各有身份标识；其负责人在2026年7月1日由甲变为乙。再假设制度规定项目由当时的负责人签批，对外联系转至负责人办公室，那么7月15日的负责人、签批人和联系办公室都应随任命改变。6月15日的负责人仍是甲，同名另一机构的负责人、甲所成立年份和历史审批的实际签字人均不受影响。这个例子把必要的依赖传播与不应发生的覆盖分开。

读图~\ref{fig:ket-update-time-scope}时，先按时间条选择事实，再沿实线检查它支持哪些结论。负责人变化影响签批人，是因为我们已经给定制度依赖，而不是因为两条文本都提到甲所。在真实系统中，没有这种规则或其他证据，就不能自动推出联系对象或审批权限也改变；相应结论应待核对。

\begin{figure*}[htbp]
\centering
% Teaching construction: appointment effective on 2026-07-01.
% Solid arrows are only the declared dependencies, not inferred associations.
\begin{tikzpicture}[
  x=1mm,y=1mm,font=\figurefont\small,
  box/.style={rectangle,draw=BookInk,fill=BookPaper,line width=0.7pt,
    align=center,inner xsep=2mm,inner ysep=2mm,minimum height=12mm},
  flow/.style={draw=BookInk,line width=0.8pt,-{Latex[length=2mm]},
    rounded corners=1mm}
]
  \node[anchor=west,text=BookInk] at (8,84) {2026年7月1日前};
  \node[anchor=west,text=BookInk] at (84,84) {2026年7月1日起};
  \draw[BookBlue,line width=2pt] (8,77) -- (80,77);
  \draw[BookTeal,line width=2pt] (80,77) -- (156,77);
  \draw[BookInk,line width=0.7pt] (80,74) -- (80,80);

  \node[box,text width=60mm,fill=FigureInput!10] (old) at (42,60)
    {旧事实：甲所负责人为甲\\有效至生效点之前};
  \node[box,text width=60mm,fill=FigureOutput!10] (new) at (122,60)
    {新事实：甲所负责人为乙\\自生效点开始有效};
  \node[box,text width=60mm] (history) at (42,34)
    {历史查询：2026年6月15日\\负责人仍为甲};
  \node[box,text width=26mm] (approve) at (101,34)
    {当前签批人\\变为乙};
  \node[box,text width=26mm] (contact) at (143,34)
    {当前联系对象\\乙的办公室};
  \draw[flow] (old.south) -- (history.north);
  \draw[flow] (new.south) -- (122,47) -| (approve.north);
  \draw[flow] (122,47) -| (contact.north);
  \node[box,text width=60mm,draw=BookMuted] at (42,10)
    {无关事实：甲所成立年份\\保持不变};
  \node[align=center,text width=70mm,text=BookMuted,
    font=\figurefont\footnotesize] at (122,10)
    {实线：给定的推理依赖\\无关事实不接依赖边};
\end{tikzpicture}

\caption{负责人任命的有效时间与依赖范围，均为教学构造。2026年7月1日分开旧、新事实的适用区间；实线按正文给定的制度连接当前签批和联系结论。历史查询沿旧事实作答，成立年份没有依赖边，因而不随任命修改。图只列已知关系，不声称已找全事实依赖。}
\label{fig:ket-update-time-scope}
\end{figure*}

按上述条件改正负责人事实及其使用，属于\term{事实编辑}（factual editing）。若要求系统在缺少有效任命来源时说明无法确认，改变的则是这种条件下的响应规则，属于\term{行为修正}（behavior correction）。前者约定何种事实应被使用，后者约定怎样响应；二者都需要明确语境。基本成功标准若随行为要求改变，第~\ref{sec:ket-transfer-output-semantics}节先确定任务，本章再实施已经明确的修正。

记录撤回提出的是另一项问题：若这条记录从未参与训练，发布结果应当怎样？\term{机器遗忘}（machine unlearning）以消除指定数据的训练影响为目标，而非只改一个答案。请求须定位记录、训练版本和派生副本，包括增强样本、特征缓存及生成标签。撤回一条含姓名的记录，不等于要求所有语境都不能出现该姓名，因为其余合法材料可能支持同一事实；屏蔽姓名同样不能证明其他输出或参数中的影响已经消除。\BookSectionRef{sec:privacy-unlearning}{卷3《基础、理论与可信性》第16章“隐私保护”第16.4节“数据撤回与模型消除”}给出了这项区别及严格保证的范围。

\subsection{保留要求与资源边界的确定}
\label{sec:ket-update-requirements}

接受雨季版本时，还要问温室任务是否继续有效，以及下一季的材料能否继续学入。\term{能力保留}（retention）比较仍服务的旧任务在更新前后的结果；前文的\term{可塑性}（plasticity）则通过后续新内容的学习过程考察继续学习的能力。旧任务分数不变，但同预算的新证据不能降低新任务风险，就不能以保留分数证明可塑性。雨季风险下降而温室风险上升，也需要报告适应的代价。\term{知识积累}（knowledge accumulation）还要求更强的对照：相关历史经验降低后续学习代价，或新经验改善仍有效的旧能力；仅仅两边都未退化不足以确认积累。

保留集合不能包含与有效请求冲突的旧要求。7月15日负责人应答乙，就应移除同条件下“必须答甲”的保护要求；6月历史查询却仍可保留甲。\emph{局部性}在这里指修改只改变请求及已核实依赖所需的响应，限制无关输入上的外溢，而不阻止必要传播。撤回要求也不能被旧行为保护抵消。因此，验收清单要分别列继续服务的任务、允许改变的条件、应吸收的新内容和禁止重新引入的记录。

表~\ref{tab:ket-update-goals}据此列出每类目标的变化对象、保留对象和参照。事实修正若同时伴随记录撤回，就要为两项主张分别提供证据，不能只选一行验收。单次编辑通过仅覆盖该请求和已测试范围；批量请求还需检查冲突，多轮维护则须复测仍有效请求、记忆覆盖和恢复条件。

\begin{table*}[htbp]
\small
\begin{tabularx}{\linewidth}{@{}p{16mm}XXXXX@{}}
\toprule
更新目标 & 期望变化 & 仍需保持的对象 & 允许的信息 & 比较参照 & 主要误判 \\
\midrule
环境适应 & 当前环境原任务风险下降 & 仍服务的原环境与类别语义 & 声明的目标输入、标签及源信息 & 原版本与不更新基线 & 把更自信当成更准确 \\
事实修正 & 指定条件下使用新事实 & 历史事实、无关关系与任务 & 有效请求、来源、时间和依赖 & 依有效条件制定的正确答案 & 只测原提示，忽略历史和关联 \\
持续吸收 & 学会当前内容并能继续学习 & 有效旧能力及可用学习容量 & 规定的经验流与记忆 & 等样本、等预算的阶段对照 & 只看遗忘，不测后续学习 \\
训练记录撤回 & 消除指定记录的训练影响 & 剩余任务及其他有效修改 & 剩余数据与允许的训练状态 & 剩余数据上的指定学习过程 & 把拒答或记录删除当成影响消除 \\
\bottomrule
\end{tabularx}
\caption{不同更新目标对应的变化、保留与验收对象。允许使用的信息在操作前声明，表中列出技术上所需的依据，不替代授权或隐私判断。多目标请求须分别满足对应行的要求，不能用其中一项成功代表整个更新。}
\label{tab:ket-update-goals}
\end{table*}

预算与保留条件一起约束可选路线。总存储限制历史证据和模型副本，更新次数限制经验的重复取用，停机上限可能排除完整重训练。账目应包含教师、优化器、生成器、编辑器、索引和验证副本，并按第~\ref{chap:ket-reuse-selection}章的原则区分准备成本与单次使用成本。只比较一次增量写入，却省略编辑器训练，回答的已不是完整路线的成本问题；核算范围须在结果出现前固定。

\section{知识更新依据的准备}
\label{sec:ket-update-evidence}

负责人修正需要任命依据，雨季适应需要能辨认失配的照片或标签，撤回需要记录到训练状态的依赖。目标不同，所需材料也不同。本节先确认访问权限，再选择历史材料、核对当前证据并划分用途。这样形成的是一份带版本与条件的证据集合；它如何被取样、赋权和转成梯度，由下一节说明。

\subsection{可用信息与访问条件的确认}
\label{sec:ket-update-access}

首先列出实际能够取得的材料，而不只写“使用历史知识”。源照片不可读时，源模型参数仍可能可用；只开放预测接口时，中间表示却未必可取。缺少源数据、目标标签和目标输入也分别排除不同操作：前两种条件仍可利用当前无标签照片，最后一种只能利用已知环境做准备。当前样本、历史表示、旧预测和修改请求均须注明来源、生成版本、适用时间及用途，才能判断它们是否提供相容的证据。

若源照片不能保留，但可读取并执行源网络，可以像SHOT那样冻结分类器，更新目标侧特征提取器。SHOT以源假设迁移（source hypothesis transfer）为思路，让新特征继续由旧分类器解释。无标签目标样本提供两个信号：减少单样本预测的不确定性并保持总体预测的多样性，以及按目标特征形成的类别中心生成伪标签；后者是模型估计的类别，作为临时监督使用而非真实鉴定。计算这些梯度仍需通过冻结分类器反向传播，所以“冻结”不是“内部不可访问”。这两种信号也不能单独识别任意新标签规律\scite{ket-liang2020-shot}。

若只允许向源服务查询概率，就不能执行上述反向传播。DINE以知识蒸馏（knowledge distillation，即让可训练学生模仿教师输出）建立目标模型，再用信息最大化继续适应。它将目标输入的源概率响应作自适应平滑，蒸馏时还约束插值输入上的预测关系，并利用单样本确定性与总体多样性。学生可以与源网络不同构，但训练需要对应输入及软概率；只得到类别编号会失去类别间的概率关系，不能原样实施完整概率匹配\scite{ket-liang2022-dine}。

用同一分类器比较三种接口，限制更容易定位。参数、结构、预处理和可微实现都可读时，可以计算中间特征和输入梯度。概率接口只给获准输入的软预测，不直接给内部梯度；有限差分可以另行估计某些变化，却有查询成本，也不等同于原生反向传播。类别接口只给决策结果，不能恢复第二候选的概率。无论哪种接口，计算旧样本预测都还需要持有该样本并获准提交；服务版本和查询预算不能由“有接口”这句话省略。

\begin{table*}[htbp]
\small
\begin{tabularx}{\linewidth}{@{}p{26mm}cccccX@{}}
\toprule
操作 & 历史原样本 & 模型参数 & 中间表示 & 源响应 & 当前标签 & 额外依赖与成本 \\
\midrule
原样本监督回放 & 必 & 学 & — & — & — & 需旧标签；保存样本并反复前后向计算 \\
历史特征回放 & — & 学 & 必 & — & — & 需旧特征及版本；只能训练特征之后的部分，除非另有原输入 \\
两域表示统计对齐 & 条 & 学 & 必 & — & — & 需两域输入或兼容的统计；编码、统计估计与批量成本 \\
重要性参数约束 & — & 必 & — & — & — & 需固定参考参数和先前估计的重要性；各需存储 \\
概率行为模仿 & — & 学 & — & 必 & — & 需输入及对应软预测；可在线查询或读取兼容缓存，另计学生训练 \\
类别行为模仿 & — & 学 & — & 必 & — & 需输入及对应硬标签，可查询或读取缓存；不足以直接匹配完整概率 \\
当前监督更新 & — & 学 & — & — & 必 & 需当前输入；标签取得和质量核对有独立成本 \\
\bottomrule
\end{tabularx}
\caption{各类更新操作直接需要的信息。“必”表示必要输入，“—”表示不必具备，“学”表示须能训练承接模型但不必读取源参数，“条”表示原样本可由兼容表示或统计替代。历史特征和源响应都须与输入及版本对应；兼容缓存可替代再次查询，不能替代输入本身。表只判断能否实施，不保证更新有效。}
\label{tab:ket-update-access}
\end{table*}

表~\ref{tab:ket-update-access}中的参数重要性是先前估计并保存的量，只有旧权重还不够。源响应也应保存查询日期、服务版本、类别顺序及能够确认的温度和预处理，未知项不能猜填；服务换版后，缓存可能已混入多个教师的行为。权限决定允许执行什么，隐私分析则考察执行后泄露什么，两者须按\BookChapterRef{chap:trust-privacy}{卷3《基础、理论与可信性》第16章“隐私保护”}分别判断。

\subsection{历史证据的选择与保存}
\label{sec:ket-update-memory}

有权限保留历史照片时，后续编码器可以重新读取其中的纹理；只存旧特征，就只能使用旧编码器已经留下的信息。在通常的压缩表示中，特征较省存储，但若细微纹理已经被丢弃，后续不能仅靠特征恢复它。历史概率保存决策及类间关系，不自动成为正确标签。生成器则把逐张保存改为按需生成，是否覆盖旧情形取决于生成质量，权重和生成时间仍需计入预算。这些载体保留的信息不同，应按下一轮需要验证的能力选择。

记忆容量不足时，先规定希望覆盖什么。随机保留面向经验流的频率，特征中心选择面向常见形态，类别配额面向必须服务的低频物种，干扰估计则面向候选更新可能损害的样本。高不确定性既可能提示决策边界，也可能来自错标或故障，不能单凭这个量决定保留价值。选择规则确定后，还要检查记忆实际漏掉哪些条件。

考虑容量$m=3$的教学记忆。依次到来的$a_1,a_2$是晴天甲类照片，$b_1$是阴影乙类，三张先填满槽位；随后到来晴天甲$a_3$和遮挡乙$b_2$。蓄水池抽样（reservoir sampling）让已见记录拥有相同的边际入选机会：第$i>m$条到来时，在$\{1,\ldots,i\}$均匀抽取$j$，仅当$j\leq m$才替换第$j$槽。因此新记录以$m/i$的概率进入，进入后等概率替换旧槽。设本次第4步抽到$j=2$、第5步抽到$j=5$，最终记忆为$\{a_1,a_3,b_1\}$。它恰好漏掉遮挡乙，但这只是一次抽样结果；每条记录在第5步后被保留的边际概率均为$3/5$。

若目标改为覆盖晴天甲、阴影乙和遮挡乙，可以在第4步用$a_3$替换同情形的$a_2$，第5步再替换冗余的$a_3$，得到$\{a_1,b_1,b_2\}$。这种规则保住稀有情形，却不再使记忆比例代表自然频率。改用预期损失增长选样时，还需根据当时模型、候选更新和可读数据估计增长，不能在以后回看已丢弃照片来补做在线选择。保存决定哪些证据仍在，训练权重决定它们贡献多少，两个决定要分开记录。

iCaRL（incremental classifier and representation learning）在类别逐步加入时，保存每类有序样本，并随类别增加缩减名额。它希望有限样本的平均特征接近该类的平均特征，因此按顺序补入能减少均值差的样本。设$\phi(\bm{x})$为归一化特征，$\bm{\mu}_k$为当前可见类$k$的平均特征，已选$r-1$张，则比较该类所有未选候选$\bm{x}$，选择使下式欧氏距离最小的一张：
\begin{equation}
\left\lVert\bm{\mu}_k-
\frac{1}{r}\left(\phi(\bm{x})+
\sum_{j=1}^{r-1}\phi(\bm{x}_{k,j})\right)\right\rVert_2
\label{eq:ket-update-herding}
\end{equation}
式中$\bm{x}_{k,j}$是此前第$j$张入选照片；此处省略平均后再归一化的实现细节，只展示逐项逼近的选择依据。新类的当前批量可以用于估计均值，旧类缩减时却只能截取已有列表前缀，不能重新访问全部旧照片。iCaRL还用保留原样本重算特征，通过比较输入到各类均值的距离预测。类别如何扩展由第~\ref{sec:ket-transfer-output-semantics}节确定，这里关注扩展后保存什么依据\scite{rebuffi2017}。

总容量固定为$m$时，$K$类平均约有$m/K$张；每类固定$m_0$张时，总量却增长为$Km_0$。照片之外，标签、时间、版本和概率缓存也占内存。特征$\phi_{\bm{\theta}^{(t)}}(\bm{x})$必须绑定生成版本，因为下一轮编码器可能改变坐标；直接混用两代特征，均值和距离便未必可比。若原输入已不在，只能保持接口不变、用证据学习兼容映射，或限制更新范围。

\subsection{当前证据与修改请求的核对}
\label{sec:ket-update-current-evidence}

两类花卉在目标特征中形成两个簇时，无标签照片能显示簇的位置，却不能决定哪簇是甲类。可靠标签可以消除这种排列歧义，但不保证未标区域的规律也已确定。最小最大熵方法MME（minimax entropy）利用已有源标签、少量目标标签及目标无标签样本，分类器以类别代表方向解释特征。有标签损失确定类别含义，无标签部分则交替让分类器增大预测熵、让编码器减小预测熵；熵衡量类别概率的不确定性，这两个方向使代表方向与目标特征共同调整。它使用已经取得的监督，不负责选择下一个标注对象\scite{ket-saito2019-mme}。

若还可以请求少量新标签，CLUE（clustering uncertainty-weighted embeddings）用不确定性加权的特征聚类安排查询。它计算目标特征与预测不确定性，聚类后选靠近中心的实例，请求标签再更新模型；聚类覆盖与不确定性共同影响选择。旧模型若对系统性错误过度自信，该权重仍可能漏选，因此覆盖抽查提供另一类证据。标注、特征提取和聚类都要计费；本章关心新增标签减少了什么失配歧义，主动获取的通用机制留给第~\ref{part:efficient-knowledge-use}部分\scite{ket-prabhu2021-clue}。

RGB与深度图要用于互相监督，还需检查配对本身。两张图成像正常，不保证对象、时刻或视角一致；拿植株甲的RGB特征监督植株乙的深度图，就是可靠记录构成了错误关系。配对须核对身份、同步、视野、配准和缺测。借第三种观测桥接时，它也须连接到同一对象或语义；共享类别名称则须核对粒度。只有教师响应时，证据仅约束查询输入上的输出，不能据此假定已获得逐像素对应。

任命请求也须逐字段核对：来源能否支持任命、机构身份是否一致、发布与生效时间有何关系，以及权限是否只覆盖某分支或代理期限。旧模型答甲是历史行为的证据，不能反驳一份有效新任命；新请求也不自动为真。身份或有效时间缺失时，应让对应结论待决，并请求能补足该字段的材料。

撤回须从记录身份追到重复副本、派生数据和训练状态。只知道文件名，不知道哪个快照用过它，就无法确定应重算的后缀；内容相似的另一条合法记录，也不能自动并入撤回集合。证据不足时，可以补监督、缩小范围、暂缓更新或保留原版本，但原版本若已违反有效撤回要求，不能直接继续服务。理论卷的不可识别性反例说明，无标签输入缺少标签规律或覆盖时，反复选择与旧模型一致的样本不会补出新答案，见\BookTheoremRef{thm:transfer-nonidentifiability}{卷3《基础、理论与可信性》第9章“迁移学习的可学习性”第9.1.2节“分布偏移与可识别性”}及其前文反例。

\subsection{更新证据与验证证据的划分}
\label{sec:ket-update-split}

一批照片若既用于改权重、调强度，又用于宣布成功，结果就可能只反映对这批照片的适配。为取得独立验收证据，本章按“更新”“选择强度及停止时机”“最终验收”划分用途，并保存分组身份、可见时间与标签到达时间。更新部分形成损失，选择部分决定配置，验收部分在候选冻结后使用。用途不同仍须检查信息依赖，不能把不同文件数当成独立样本数。

例如，教学上取互不重叠的植株采集组$g_1,g_2,g_3$，每组连拍整体分配，得到以下用途记录：
\[
\begin{array}{ccl}
g_1&\longrightarrow&\text{更新集：计算梯度}\\
g_2&\longrightarrow&\text{选择集：选约束强度与停止时机}\\
g_3&\longrightarrow&\text{验收集：候选冻结后评估}
\end{array}
\]
组号只说明用途划分，不规定比例。实际划分还须覆盖待验证环境和类别；同组近重复照片不能因文件不同就分到训练与独立验收两侧。

事实改写也存在这种依赖。编辑训练与泛化验证可以分别保留表达，但模板、实体和事实族的划分须事先约定；已经帮助选编辑器的表达不能再充当独立成功证据。相邻时间窗口可能共享环境状态及延迟标签，应按实际服务时序切分，避免未来反馈参与过去决策。

连续更新时，用途须随数据流维护。星期一照片无标签，系统先预测并作获准的无标签适应；星期三鉴定返回，只能给星期一已经保存的预测计分，或用于以后的监督更新，不能用新模型重答后替换原结果。延迟标签后来若转入训练，其独立验证身份也随之结束，后续版本要另留验收证据。

测试输入也有使用边界。整批共同估计归一化统计，与逐样本只读过去状态不同；反复遍历完整输入集，又与单遍流不同。即便不读测试标签，提前联合未来输入仍改变信息条件。因此须记录批量、顺序、跨样本状态、重置边界和计分时刻，后文才能把实际计算与所称协议对应起来。

\section{知识的更新}
\label{sec:ket-update-operations}

现在已有明确目标和可用证据，需要决定把变化写入哪里。固定结构内改权重、偏置或校准量，是参数更新；增减组件和连接、调整允许的路由及容量，是结构操作；修改可调用记录及有效条件，是外部更新。三者按修改对象区分，可以组合。例如，为深度图增加编码器改变结构，训练编码器改变参数；写入任命文件改变记录，训练时间条件识别器又改变参数。

操作规模取决于需要重估的关系。广泛成像失配需要覆盖这些条件的经验；高频变化且带时间的事实，可以先考虑可追溯的外部记录；能定位并检验影响范围的请求，也可考虑局部参数编辑。完整重训练重估范围较大，但须处理训练集之外的有效编辑；局部方法减少了改动范围，却不自动满足关联一致性或撤回标准。成本比较应覆盖各路线所需的完整过程，而不预设局部方法总是更便宜。

\subsection{模型参数的更新}
\label{sec:ket-update-parameters}

即使结构固定，也还不能从数据直接得到唯一的参数增量。样本及权重决定经验的贡献，目标和约束规定希望改变的关系，可训练范围及起点限制可达解，优化过程决定怎样产生候选版本。以下按这些选择的依赖展开。回放样本既可提供标签损失，也可提供历史行为目标；对齐损失也可作用于局部层或全模型。定向编辑和共享基更新横跨多个选择，讨论范围时将连同它们的计算解释，而不把这些机制当成互斥类别。

\subsubsection{训练经验的取用与加权}
\label{sec:ket-update-weighting}

原训练集的平均损失反映源频率，未必反映雨季频率。输入输出含义相同且目标情形被源覆盖时，我们可以改变旧样本的贡献，而不先构造新照片。\term{重要性加权}（importance weighting）以目标与源的概率比赋权：目标相对更常见的情形贡献更多，源样本上的加权平均由此对应目标风险。这种对应首先是分布恒等式，还需估计权重才能用于训练。

为求出这项权重，设联合密度$p_S,p_T$相对于同一基准测度定义，损失可积，并要求$p_T(\bm{x},y)>0$时$p_S(\bm{x},y)>0$。将$R_T(f_{\bm\theta})$简记为$R_T(\bm\theta)$，对离散标签先求和，再对输入积分；在源覆盖区域乘除$p_S$，就可把目标风险改写为源分布上的期望：
\begin{align}
R_T(\bm{\theta})
&=\sum_y\int \ell(f_{\bm{\theta}}(\bm{x}),y)
       p_T(\bm{x},y)\,\mathrm{d}\bm{x}\notag\\
&=\sum_y\int \ell(f_{\bm{\theta}}(\bm{x}),y)
       \frac{p_T(\bm{x},y)}{p_S(\bm{x},y)}
       p_S(\bm{x},y)\,\mathrm{d}\bm{x}\notag\\
&=\mathbb E_{(\bm{x},y)\sim\mathcal P_S}
  [w(\bm{x},y)\ell(f_{\bm{\theta}}(\bm{x}),y)].
\label{eq:ket-update-importance}
\end{align}
密度比$w=p_T/p_S$将源概率转换成目标概率，源零概率且目标也为零处的取值不影响期望。连续标签则将外层求和改为积分。无标签目标输入通常不能估计完整联合密度比；偏移假设的用途，是消去不可见的标签项，而不是由等式直接保证可实施。

协变量偏移保留$p(y\mid\bm{x})$，因此比值中的条件项相消，只留下可由两域输入估计的密度比。用$n_S$个带标签源样本及估计权重$\hat w$，形成经验目标：
\begin{equation}
\begin{aligned}
w(\bm{x},y)
&=\frac{p_T(\bm{x})p_T(y\mid\bm{x})}
       {p_S(\bm{x})p_S(y\mid\bm{x})}
=\frac{p_T(\bm{x})}{p_S(\bm{x})},\\
\hat R_w(\bm{\theta})
&=\frac{1}{n_S}\sum_{i=1}^{n_S}
\hat w(\bm{x}_{S,i})\ell(f_{\bm{\theta}}(\bm{x}_{S,i}),y_{S,i}).
\end{aligned}
\label{eq:ket-update-covariate}
\end{equation}
下标$S,i$标记源样本，帽号说明权重只是估计值。除直接估计密度比，还可训练一个区分源、目标输入的域分类器。若它在混合数据上估计目标域概率$d(\bm{x})$，两域采样先验为非零的$\pi_S,\pi_T$，贝叶斯公式给出
\[
\frac{p_T(\bm{x})}{p_S(\bm{x})}
=\frac{d(\bm{x})}{1-d(\bm{x})}\frac{\pi_S}{\pi_T}.
\]
这个概率相对于\emph{训练域分类器时}的混合分布定义，所以平衡采样与自然频率采样的先验系数不同。公式中的赔率要求$d(\bm{x})<1$；有限样本的饱和输出还需数值处理。权重估计和相关配置应在验收前冻结，并检查误差，不能把分类分数当成已知精确比值。

教学上把输入分成不重叠的晴天与阴影情形，源比例为$(0.8,0.2)$，目标为$(0.4,0.6)$，情形内的输入与条件标签规律都保持。固定识别器的两项平均损失设为$(0.1,0.3)$，源风险为$0.8\times0.1+0.2\times0.3=0.14$。权重$(0.5,3)$给出贡献$0.04$和$0.18$，总量$0.22$，恰等于目标比例下的风险。变化落在每条损失的权重，而不是对最终准确率任意乘比例。

阴影权重较大，也让少数阴影样本的误差更影响估计。令$Z=w\ell$具有有限方差，独立同分布源样本的均值方差为$\operatorname{Var}_S(Z)/n_S$。为限制单条影响，可用阈值$c>0$取$w_c=\min(w,c)$；若$0\leq\ell\leq L$，则截断损失与目标风险之差满足
\[
0\le R_T-\mathbb E_S[w_c\ell]
=\mathbb E_S[(w-c)_+\ell]
\le L\,\mathbb E_S[(w-c)_+].
\]
其中$(u)_+=\max(u,0)$。上界用被截去的权重质量控制偏差，说明限制影响不是无代价的处理。若再归一化截断权重，偏差表达也会改变，不能沿用无偏结论。目标区域若根本没有源覆盖，则前面的恒等式已不适用，增大有限权重也不能创造缺失标签。

若保持的是每类输入分布，就按类别而非拍摄情形校正。记先验$\pi_{S,k}=p_S(y=k)$、$\pi_{T,k}=p_T(y=k)$，要求目标出现的类在源中非零。标签偏移使类条件密度相消，得到
\begin{equation}
w(\bm{x},y=k)
=\frac{p_T(\bm{x}\mid k)\pi_{T,k}}
       {p_S(\bm{x}\mid k)\pi_{S,k}}
=\frac{\pi_{T,k}}{\pi_{S,k}}.
\label{eq:ket-update-label-weight}
\end{equation}
同一类别内各样本因此共享一个权重。另设两类先验从$(0.8,0.2)$变成$(0.4,0.6)$，类内分布及教学损失$(0.1,0.3)$保持，贡献也得到$(0.04,0.18)$。重复这组数值是为了检验加权对象：前例按晴天或阴影，本例按物种，算术相同不能使假设相同。

还不知道目标先验时，可以利用固定预测器对各真类的响应差异。令$\hat y=h(\bm{x})$，定义列按真实类别归一化的混淆矩阵，以及无标签目标输入的预测频率：
\[
\bm C_{j,k}=\mathbb P_S(\hat y=j\mid y=k),
\qquad q_{T,j}=\mathbb P_T(\hat y=j).
\]
固定$h$且类条件分布保持，就有$\mathbb P_T(\hat y=j\mid y=k)=\bm C_{j,k}$。按真类分解目标预测频率得到线性关系。有限样本时，在非负且和为一的类别比例中选择最能拟合该关系的估计：
\begin{equation}
\bm q_T=\bm C\bm\pi_T,
\qquad
\hat{\bm\pi}_T
=\argmin_{\bm\pi\ge0,\ \sum_k\pi_k=1}
\lVert\hat{\bm C}\bm\pi-\hat{\bm q}_T\rVert_2^2.
\label{eq:ket-update-bbse}
\end{equation}
这里用平方残差衡量预测频率的拟合误差，约束只保证估计是合法概率。$\hat{\bm C}$须由与预测器拟合适当隔离的源标签估计，$\hat{\bm q}_T$由许可的目标输入估计。方阵$\bm C$可逆时，总体有$\bm\pi_T=\bm C^{-1}\bm q_T$，有限样本直接求逆却可能产生负数；接近奇异时，合法概率约束也不能消除误差放大。黑盒偏移估计BBSE（black box shift estimation）的原始写法可用联合混淆矩阵求先验比；这里求的是列归一化形式下的目标先验，未知量须区分\scite{ket-lipton2018-labelshift}。

用一个两类教学例子复算可识别条件，取
\[
\bm C=
\begin{pmatrix}0.9&0.2\\0.1&0.8\end{pmatrix},
\qquad
\bm q_T=
\begin{pmatrix}0.48\\0.52\end{pmatrix}.
\]
第一行给出$0.48=0.9\pi_{T,1}+0.2(1-\pi_{T,1})$，故先验为$(0.4,0.6)$；$\det(\bm C)=0.7$保证本例可逆。所需条件不是预测完全正确，而是真类产生可区分的预测分布。若两列相同，改变类比例不会改变$\bm q_T$，再多无标签预测也无法辨认比例。

得到先验后，可以重新加权训练，也可以只校正预测。先对真实源后验应用贝叶斯公式，在标签偏移下得到
\begin{equation}
p_T(k\mid\bm{x})
=\frac{p_S(k\mid\bm{x})\,\pi_{T,k}/\pi_{S,k}}
       {\sum_j p_S(j\mid\bm{x})\,\pi_{T,j}/\pi_{S,j}}.
\label{eq:ket-update-prior-correction}
\end{equation}
分母把按先验比缩放的分数归一化。若分类器的logit（进入softmax前的类别分数）$a_k(\bm{x})$给出源后验，可加偏置$\log\pi_{T,k}-\log\pi_{S,k}$，不改编码器和头的原权重；零目标先验对应排除该类的极限处理。源后验$(0.6,0.4)$在本例变为未归一化$(0.3,1.2)$，归一化后为$(0.2,0.8)$。实际网络只近似后验时，校正也只是近似。BBSE利用硬预测混淆关系，不要求概率校准；后验校正却依赖概率质量。类条件规律变化后，两者都不能给出一般适应保证。

校正当前频率之外，旧任务也需要继续提供训练信号。\term{经验回放}（experience replay）将已保存样本重新送入训练，让旧情形仍能影响梯度。设非空的当前集$S_{\mathrm{new}}$与历史记忆$M$已经选定，分别求平均后，以旧经验份额$\rho$组合：
\begin{equation}
L_{\mathrm{mix}}(\bm{\theta})
=(1-\rho)\frac{1}{|S_{\mathrm{new}}|}
\sum_{(\bm{x},y)\in S_{\mathrm{new}}}\ell(f_{\bm{\theta}}(\bm{x}),y)
+\rho\frac{1}{|M|}\sum_{(\bm{x},y)\in M}
\ell(f_{\bm{\theta}}(\bm{x}),y),
\label{eq:ket-update-replay}
\end{equation}
系数$0\leq\rho\leq1$由保留与吸收要求选择，不由库中样本数自动决定。每步分别均匀抽取新、旧批量，求各自均值再组合，才能对应这项目标；拼接后统一平均，则由批量比例隐式赋权。小记忆回放研究采用当前批量、历史批量共同更新并维护记忆的流程，为比较复杂保留方法提供直接混合基线；本文不据此断言它在所有任务上优于其他方法\scite{ket-chaudhry2019-replay}。

教学上取100张新照片、20张记忆，每个新批量10张。十步每步回放10张，共有100次新样本和100次旧样本处理；每两个新批量才回放10张，则旧处理降为50次。记忆容量相同，训练成本和旧信号份额并不相同。间隔回放改为每次20张虽恢复100次处理，顺序、批量统计和中途漂移仍有差异。因此频率比较必须记录处理次数及损失系数，而不只记库容量。

采样概率与损失系数还可分别改变。固定参数下，若目标为$\sum_i a_i\ell_i$，$a_i\geq0$且和为一，而按概率$q_i$抽样，则所有$a_i>0$处须有$q_i>0$，单次贡献$a_i\ell_i/q_i$的期望才等于目标。不校正时，平均估计的是采样分布下的另一目标。主动提高稀有物种权重可以服务覆盖或公平要求，但应将这种取舍说清，不能仍称它为自然频率风险。

旧样本也能携带历史输出，而不只携带标签。暗经验回放DER（dark experience replay）保存输入及历史logit向量$\bm a_i^{\mathrm{old}}$，希望当前logit别偏离这些记录；平方差对大偏离给予更强惩罚。与新任务损失组合后，目标为
\begin{equation}
L_{\mathrm{DER}}(\bm{\theta})
=L_{\mathrm{new}}(\bm{\theta})
+\lambda_{\mathrm{logit}}\,
\mathbb E_{(\bm{x}_i,\bm a_i^{\mathrm{old}})\sim M}
\lVert\bm a_{\bm{\theta}}(\bm{x}_i)-\bm a_i^{\mathrm{old}}\rVert_2^2.
\label{eq:ket-update-der}
\end{equation}
其中$L_{\mathrm{new}}$是当前任务目标，$\lambda_{\mathrm{logit}}\geq0$选择历史分数约束的强度，期望按记忆抽样计算。缓存来自不同历史时刻，并非都来自任务终点，故类别坐标和版本必须对应。输入回放确定覆盖范围，logit匹配确定要保留的输出，两者共同起作用。DER++另加旧真实标签损失，以减少固定错误历史输出的风险；这项监督不属于原始DER的必选项\scite{ket-buzzega2020-der}。

生成器产生旧输入，再提供标签或教师响应，也可进入回放及DER目标，但平均量此时对应生成分布。若它漏掉罕见阴影物种，多生成常见照片不会补足覆盖；教师错误又可能被反复模仿。生成器继续训练也可能丢失旧情形，所以须同时检查生成质量、覆盖、教师版本和成本，不能把生成样本等同于原历史证据。

把回放用于持续预训练时，经验单位从照片换成词元序列，即模型读入的文本片段；指令学习还需把指令与响应配对。文档等比例不代表词元等比例，长文档与重复指令会改变梯度份额。混合新旧材料、先学新域再回放，或周期插入通用材料，不仅顺序不同，总处理量也可能不同。比较时需记录词元和样本权重，再检查领域风险、旧能力、后续学习曲线与累计成本。

未来目标输入还不可用时，训练变化须来自已知条件。\term{域随机化}（domain randomization）在仿真或可控生成中改变光照、纹理和视角，让同一任务经历不同观测；Tobin等以仿真物体定位展示这一路线。花卉应用还须限定哪些变化保留物种标签，不能任意改变判别所需的花瓣结构。扩展渲染条件增加的是已定义的观测范围，不等于覆盖全部现实环境\scite{ket-tobin2017-randomization}。

更一般地，\term{领域泛化}（domain generalization）只用已知训练环境准备和选择下一版本，不使用待部署域的数据。若服务要求不能只顾环境平均值，可以在这些已声明环境$e$上最小化$\max_e\hat R_e(\bm{\theta})$：最大值取经验风险最高的已知环境，防止其被平均掩盖。它并未对所有未来域取最大值，有限样本噪声也可能支配这项目标，需要验证和正则化\scite{sagawa2020groupdro}。

另一种选择是不变风险最小化（invariant risk minimization）：在各训练环境上，要求共享表示后的同一预测器满足各自的最优性条件，实际优化再作近似。它约束预测关系，而不只使输入统计相同。若所有已知域都让背景颜色同样预测物种，数据仍无法辨认这种关联会在未来失效。环境划分及连接未来环境的假设因而决定解释范围，理论边界见\BookSectionRef{sec:domain-generalization}{卷3《基础、理论与可信性》第9章“迁移学习的可学习性”第9.3节“领域泛化的可学习性”}\scite{arjovsky2019irm}。

\subsubsection{更新目标与保留约束的构造}
\label{sec:ket-update-objective}

样本贡献确定后，还需规定要学什么关系。若雨季照片保留物种形态，编码器却把阴影当作类别线索，就应调整表示；温室仍需服务时，还要保护有效旧预测。将模型分为编码器$\phi_{\bm{\omega}}$与预测头$h_{\bm{\psi}}$，写$f_{\bm{\theta}}=h_{\bm{\psi}}\circ\phi_{\bm{\omega}}$、$\bm{\theta}=(\bm{\omega},\bm{\psi})$。以下先构造对应与对齐目标，再说明如何用历史参数、输出和表示限制无谓改动。

同一特征的波动及特征间的共同变化若因域而异，可以先比较协方差。Deep CORAL以深度相关对齐（deep correlation alignment）将这种比较加入编码器学习。对域$D\in\{S,T\}$的$n_D>1$个$d$维表示，先求均值，再求中心化协方差：
\[
\bar{\bm z}_D=\frac1{n_D}\sum_i\bm z_{D,i},\qquad
\bm C_D=\frac1{n_D-1}\sum_i
(\bm z_{D,i}-\bar{\bm z}_D)(\bm z_{D,i}-\bar{\bm z}_D)^\top.
\]
两域使用相同特征坐标时，以协方差矩阵各元素差的平方和衡量二阶统计失配，得到
\begin{equation}
L_{\mathrm{CORAL}}=\frac1{4d^2}\lVert\bm C_S-\bm C_T\rVert_{\mathrm F}^2.
\label{eq:ket-update-coral}
\end{equation}
Frobenius范数的平方$\lVert\cdot\rVert_{\mathrm F}^2$累计元素平方，$1/(4d^2)$控制维数带来的尺度。训练时源标签提供任务损失，两域表示提供对齐损失，共同更新编码器；部署只保留任务模型。协方差相同不代表均值或完整分布相同。若用历史协方差缓存代替源输入，还须证明旧、当前坐标兼容，矩阵形状相同不能保证比较有意义\scite{ket-sun2016-coral}。

不预先固定二阶距离时，可以让可训练判别器寻找域差异。域对抗神经网络DANN（domain-adversarial neural network）让判别器辨认来源，同时让编码器削弱这种辨认能力。令$d_{\bm{\nu}}(\bm z)$输出目标域概率，用两域输入的负对数概率评价其错误：
\[
L_{\mathrm{dom}}
=-\mathbb E_S\log(1-d_{\bm{\nu}}(\phi_{\bm{\omega}}(\bm x)))
 -\mathbb E_T\log d_{\bm{\nu}}(\phi_{\bm{\omega}}(\bm x)).
\]
两项期望分别对源、目标输入取值；判别器降低$L_{\mathrm{dom}}$，编码器则降低$L_{\mathrm{task}}-\lambda_{\mathrm{dom}}L_{\mathrm{dom}}$，其中非负系数控制域信号强度。梯度反转层前向传递表示，反向对编码器的域梯度乘负系数，实现相反方向；任务头仍按标签学习。部署一般不调用域判别器。它比固定协方差度量多了一项可学习的差异检测，但检测失败也可能来自容量不足或未训练好，不能直接当作分布相同的证据\scite{ganin2016}。

这些方法对齐边际统计，却还没有确定目标类与源类怎样对应。设两域各有两个等大簇，源左为甲、右为乙，目标编码器却将甲放右、乙放左。边际分布完全一致，旧头仍把全部目标分错。常量编码也能消除域差，却抹掉分类信息。源任务损失可排除部分退化编码，不能单独辨认全部目标标签排列；理论卷在\BookTheoremRef{thm:transfer-nonidentifiability}{卷3《基础、理论与可信性》第9章“迁移学习的可学习性”第9.1.2节“分布偏移与可识别性”}之前给出了相同可观察输入对应相反标签的反例。

条件域对抗网络CDAN（conditional domain adversarial network）进一步把类别预测作为对齐条件。判别器读取$\bm z\otimes\bm p$，其中$\bm p$是类别概率，外积将特征分量与每类概率相乘，从而比较特征和预测的联合关系。任务监督与对抗仍共同训练模型，熵加权可降低不确定预测的贡献；特征或类别较多时，可用随机多线性近似减小输入维数。条件来自模型预测，而不是真实目标标签，所以系统性错配仍可能被强化\scite{ket-long2017-cdan}。

若希望把有限样本间的对应显式保存，可以用最优传输（optimal transport）分配质量。设源有$n_S$点、目标有$n_T$点，非负权重$\bm a,\bm b$各自和为一，代价矩阵$\bm D\in\mathbb R^{n_S\times n_T}$由可比较表示间的距离等规则确定。允许的耦合须保持两侧质量，再选择总代价与熵正则之和最小的一个：
\begin{equation}
\min_{\bm\Gamma\ge0}
\sum_{i,j}\bm\Gamma_{i,j}\bm D_{i,j}
+\varepsilon\sum_{i,j}\bm\Gamma_{i,j}\log\bm\Gamma_{i,j},
\quad
\bm\Gamma\bm 1=\bm a,\quad
\bm\Gamma^\top\bm1=\bm b.
\label{eq:ket-update-transport}
\end{equation}
$\bm\Gamma_{i,j}$是从源点$i$分给目标点$j$的质量，行、列和约束防止凭空增减质量，$\bm1$取相应维数，$\varepsilon\geq0$控制分配的分散程度，约定$0\log0=0$。平方距离下，$a_i>0$点的重心映射为$\tilde{\bm z}_{S,i}=\sum_j\bm\Gamma_{i,j}\bm z_{T,j}/a_i$，可在这些映射后的带标签点上学目标坐标中的预测器，也可将代价用于表示训练。类别或邻域约束还可限制混合。新输入须使用已经学得的预测器或外推映射，有限耦合本身不是编码器；几何代价若不反映物种语义，最小代价也可能选错对应\scite{ket-courty2015-transport}。

对齐须同时维护判别信息。花色既可能反映相机响应，也可能区分物种，能预测域身份不意味着应全部消除。对齐强度要用许可的选择证据确定；有目标标签时检查类内、类间与任务误差，无标签时只能报告统计代理和所采用假设，不能宣称已恢复类别关系。

RGB换成深度图时，不再直接比较原输入坐标。跨模态蒸馏（cross-modal distillation）用同一对象的旧观测为新观测提供学习目标。设$n>0$对样本为$(\bm x_{A,i},\bm x_{B,i})$，$A$为旧模态，$B$为新模态，冻结旧编码器。若目标是让学生读出教师特征，可以最小化配对特征的平均平方误差：
\begin{equation}
L_{\mathrm{pair}}
=\frac1n\sum_i
\left\lVert g_{\bm\eta}(\phi_{B,\bm\omega}(\bm x_{B,i}))
-\phi_{A,\mathrm{old}}(\bm x_{A,i})\right\rVert_2^2.
\label{eq:ket-update-crossmodal}
\end{equation}
学生编码器$\phi_{B,\bm\omega}$与映射$g_{\bm\eta}$可学习，后者输出维数与教师特征一致，原输入维数则不必相同。训练读取配对旧观测和教师，部署只需新观测、学生及映射。原始研究利用配对RGB与深度或光流传递监督；这里采用平方特征误差解释一种匹配形式，关键是对象对应，不是两批记录共享一个宽泛类别\scite{ket-gupta2015-crossmodal}。

广义蒸馏（generalized distillation）还允许教师读取部署时不可得的特权信息（privileged information），学生只读自身可用输入，同时学习标签与教师软预测。额外训练信息能提供更细目标，却不能消除新输入的不可识别性。若两物种仅花色不同，深度图分布完全相同，确定映射就不能逐张恢复花色；有限二阶矩及不受限函数条件下，平方误差的总体最优输出是给定深度的教师特征条件均值，不是每张真实RGB特征。共享语义还须由配对、标签或桥接证据建立，同为$d$维不够。后文第~\ref{sec:ket-integration-student-training}节展开通用蒸馏流程，第~\ref{part:efficient-knowledge-use}部分承担多模态信息获取的内容边界\scite{ket-lopezpaz2016generalized}。

跨语言时，对应对象从同一植株的观测换成同一含义的表达。平行文本与实体标识可以连接词序、词形，却须核对专名、日期和否定范围。共享编码器若已保留所需语义，只是语言材料比例变了，可以先检验混合权重；关键表达未被覆盖时，则需补对应监督。若输出规范和成功标准也变，先由第~\ref{sec:ket-transfer-output-semantics}节确定任务，不能只校正输入频率。

前述加权和匹配思想不只用于图像，但换任务要重建条件。回归目标连续，须保留单位及条件数值关系，离散类先验公式不能原样套用；厘米与米混用应先修接口。时序输入依赖历史窗口，未来标签不能决定当前对齐，重排也可能破坏时间关系。图节点表示依赖邻居、边类型和传播，节点或结构更新会影响其他节点，跨图对应须核对实体及关系。可复用的是构造目标的思路，不是分类实验对这些新范围的自动保证。

新规律获得证据后，还要限制它对有效旧能力的影响。让所有参数等距离靠近$\bm\theta^\star$，没有区分旧损失敏感与不敏感的方向。\term{弹性权重巩固}（elastic weight consolidation, EWC）依据局部敏感程度加权偏移，使重要方向受到较强惩罚。历史后验的局部展开给出了这项二次约束的来源。

设历史集$S_{\mathrm{old}}$与新集$S_{\mathrm{new}}$在给定参数后条件独立，负对数后验可分解为
\[
-\log p(\bm\theta\mid S_{\mathrm{old}},S_{\mathrm{new}})
=-\log p(S_{\mathrm{new}}\mid\bm\theta)
-\log p(\bm\theta\mid S_{\mathrm{old}})+\mathrm{const}.
\]
第一项需要新数据，第二项保存历史约束。若$\bm\theta^\star$是可微历史后验的内部局部极大点，一阶项为零；以历史负对数后验的Hessian$\bm H_{\mathrm{old}}$作局部二阶展开，得到
\[
-\log p(\bm\theta\mid S_{\mathrm{old}})
\approx\mathrm{const}
+\tfrac12(\bm\theta-\bm\theta^\star)^\top
\bm H_{\mathrm{old}}(\bm\theta-\bm\theta^\star).
\]
这个展开仍含完整曲率矩阵。适当正则条件下，模型分布的负对数似然期望曲率对应Fisher信息；EWC以其对角近似形成较易保存的敏感性权重：
\begin{align}
F_j&=\mathbb E_{\bm x\sim\mathcal P_{\mathrm{old}},\
y\sim p_{\bm\theta^\star}(\cdot\mid\bm x)}
\left[
\left(\frac{\partial}{\partial\theta_j}
\log p_{\bm\theta^\star}(y\mid\bm x)\right)^2
\right],\notag\\
L_{\mathrm{EWC}}(\bm\theta)
&=L_{\mathrm{new}}(\bm\theta)
+\frac{\lambda_{\mathrm E}}2
\sum_j\hat F_j(\theta_j-\theta_j^\star)^2.
\label{eq:ket-update-ewc}
\end{align}
这里$\mathcal P_{\mathrm{old}}$仅表示历史输入分布，$y$从参考模型条件概率抽取；$\hat F_j$估计第$j$个参数的平方评分梯度。非负$\lambda_{\mathrm E}$吸收样本与损失尺度及近似系数。观测标签替代模型采样标签得到的经验Fisher，不无条件等于总体Fisher或样本Hessian。对角化又舍弃参数相关方向，因此二次罚项只近似历史局部约束，不是完整后验或无遗忘保证\scite{kirkpatrick2017ewc}。

新阶段前固定$\bm\theta^\star$和$\hat F_j$后，即便不再保存旧照片，仍可计算罚项，但历史证据已经用于估计重要性，并留下两份参数规模向量。其梯度$\lambda_{\mathrm E}\hat F_j(\theta_j-\theta_j^\star)$拉回偏离；参考若每步随动，就不再保护同一历史状态。原始任务式EWC利用任务终点巩固，多任务如何合并近似须声明，无边界流也不能默认提供分段时刻。

突触智能SI（synaptic intelligence）改用训练轨迹估计贡献，避免另在终点逐项测曲率。对一个已结束学习区间，$g_j^{(s)}$表示任务梯度，$\Delta\theta_j^{(s)}$表示实际位移；局部损失减少量由二者负乘积近似。以下累计贡献，再按净位移归一化，并明确截去负估计：
\[
\omega_j=-\sum_s g_j^{(s)}\Delta\theta_j^{(s)},\qquad
\Omega_j\leftarrow\Omega_j+
\frac{\max(\omega_j,0)}{(\theta_{j,\mathrm{end}}-\theta_{j,\mathrm{start}})^2+\xi},
\]
$\xi>0$避免零净位移的除零，$\Omega_j$从零累计。$\omega_j$是沿路径的贡献近似，受随机梯度、优化器和步长影响；负值截断是此处声明的处理，而非轨迹必然非负的结论。区间结束后固定重要性及参考，以加权平方偏离约束下一阶段。SI不需为该累计另存全部旧样本，却需轨迹状态和巩固时刻。它与EWC测量不同近似下的重要性，都不能将一个事实唯一定位到某个参数\scite{ket-zenke2017-si}。

参数距离不能直接规定答案，若关心旧回答，可以匹配有效输入上的概率。取非空保留集合$U$，其元素须获准查询且不与当前修改冲突，冻结旧模型。温度$T>0$共同缩放两侧logit，形成
\[
q_{i,k}=\frac{\exp(a^{\mathrm{old}}_{i,k}/T)}
{\sum_j\exp(a^{\mathrm{old}}_{i,j}/T)},\qquad
p_{i,k}(\bm\theta)=\frac{\exp(a_{i,k}(\bm\theta)/T)}
{\sum_j\exp(a_{i,j}(\bm\theta)/T)}.
\]
温度越高通常使分布越平缓，但不会使旧预测变成真值。两侧类别顺序与输出范围须一致，旧概率$q$固定。用相对熵衡量当前分布偏离旧分布的程度，再对$U$平均，得到
\begin{align}
L_{\mathrm{beh}}(\bm\theta)
&=\frac{T^2}{|U|}\sum_{i\in U}\sum_k
q_{i,k}\log\frac{q_{i,k}}{p_{i,k}(\bm\theta)}\notag\\
&=\mathrm{const}
-\frac{T^2}{|U|}\sum_{i\in U}\sum_k q_{i,k}\log p_{i,k}(\bm\theta).
\label{eq:ket-update-behavior}
\end{align}
相对熵在概率一致时为零；因$q$固定，第二行与软标签交叉熵仅差常数。对单个学生logit求导得$\frac{T}{|U|}(p_{i,k}-q_{i,k})$，梯度下降因而向旧概率调整，仍须保留批平均因子。$T^2$补偿一部分温度引起的尺度变化，不保证各温度的目标与训练轨迹等价。

缓存若保存精确完整softmax概率且知道温度$T_0>0$，可取各概率的$T_0/T_1$次幂再归一化，恢复同一logit在$T_1>0$下的分布，因为共同加性常数会消去。硬类别、截断概率、舍入误差或未知后处理破坏这项恢复条件，不能由任意缓存宣称精确换温度。

无遗忘学习LwF（learning without forgetting）以当前任务输入查询旧模型的旧任务输出，把行为匹配与当前监督结合。它不要求回放旧原样本，却仍需教师或兼容缓存；当前输入若不覆盖旧有效区域，那里便缺少保护证据。部署保留新网络及任务头，不需持续调用教师。DER明确回放历史输入与轨迹分数，LwF典型地查询当前输入，二者的区别首先在约束覆盖什么，而不只在是否称为蒸馏\scite{ket-li2016-lwf}。

没有即时标签时，也可保护旧表示中可读的信息。持续自监督学习CaSSLe（continual self-supervised learning）冻结历史编码器，让当前特征经可学习预测器匹配旧特征，同时学习当前自监督目标，即以输入自身构造监督。匹配沿用具体方法的相似性或对比损失，不一律是平方误差。预测器允许当前坐标变化而保持旧特征可读；旧编码器与当前输入提供信号，缺少旧区域覆盖的限制仍在。后来用标签训练线性探测器，只说明表示可读性；没有当时的任务头，不能把探测成绩算成即时回答能力\scite{ket-fini2021-cassle}。

为检查各约束是否指向同一候选参数，可将非对抗损失与参数罚项写成一个分析用组合：
\begin{equation}
\begin{split}
L(\bm\theta)={}&L_{\mathrm{new}}(\bm\theta)
+\lambda_{\mathrm R}L_{\mathrm{replay}}(\bm\theta)
+\lambda_{\mathrm A}L_{\mathrm{align}}(\bm\theta)\\
&+\lambda_{\mathrm B}L_{\mathrm{beh}}(\bm\theta)
+\frac{\lambda_{\mathrm P}}2
\sum_j\Omega_j(\theta_j-\theta_j^\star)^2.
\end{split}
\label{eq:ket-update-joint-objective}
\end{equation}
$L_{\mathrm{replay}}$是记忆任务损失，$\lambda_{\mathrm R}$相对系数为1的新损失赋权，与归一化份额$\rho$不同；$L_{\mathrm{align}}$依赖两域或配对证据，$L_{\mathrm{beh}}$依赖有效保留输入，$\Omega_j$取已解释的重要性。各$\lambda$非负，由选择证据确定，不需要全部启用。对抗还须单独更新判别器，不能让所有参数共同下降此式。即使复用同一批输入，统计、行为和参数罚项也分别约束表示关系、输出与位移，不能相互替代。

目标相加前先排除语义冲突。7月15日负责人应改乙，却若在同输入上严格保护甲，就没有同时满足两条要求的预测器；提高两系数只改变软折中，不修复矛盾。先移除已失效的旧保护条件，仍保留6月历史查询与无关实体，再对允许折中的任务选强度。图~\ref{fig:ket-update-objective}将固定教师、重要性与参考参数画在约束侧，当前量画在计算侧，可据此检查参照是否随优化误动。

\begin{figure*}[htbp]
\centering
% All losses refer to the same current parameter vector.
% Fixed teacher and parameter references must not move with the student.
\begin{tikzpicture}[
  x=1mm,y=1mm,font=\figurefont\small,
  box/.style={rectangle,draw=BookInk,fill=BookPaper,line width=0.7pt,
    align=center,inner xsep=2mm,inner ysep=2mm,minimum height=12mm},
  flow/.style={draw=BookInk,line width=0.8pt,-{Latex[length=2mm]},
    rounded corners=1mm},
  constraint/.style={flow,dashed,draw=BookBlue},
  source/.style={box,text width=37mm},
  computed/.style={box,text width=32mm,fill=FigureOutput!10},
  loss/.style={box,fill=FigureLoss!10,text width=24mm}
]
  \node[anchor=west,text=BookInk] at (1,98)
    {各行共用当前模型$f_{\bm\theta}$，历史参照在本轮固定};
  \node[source,fill=FigureInput!10] (data) at (22,80)
    {当前与回放批量\\预测输入及有效标签};
  \node[source,fill=FigureInput!10] (pairs) at (22,56)
    {两域或配对预测输入\\明确对应依据};
  \node[source,fill=FigureModel!10] (teacher) at (22,32)
    {固定旧模型与预测请求$U$\\提供历史预测$\bm q$};
  \node[source,fill=FigureModel!10] (reference) at (22,8)
    {固定参考$\bm\theta^\star$\\固定重要性$\bm\Omega$};
  \node[computed] (prediction) at (70,80)
    {当前预测\\$f_{\bm\theta}(\bm x)$};
  \node[computed] (features) at (70,56)
    {当前表示\\$\phi_{\bm\omega}(\bm x)$};
  \node[computed] (behavior-pair) at (70,32)
    {当前$\bm p_{\bm\theta}(U)$\\与固定$\bm q$比较};
  \node[computed] (difference) at (70,8)
    {当前模型参数偏离\\$\bm\theta-\bm\theta^\star$};
  \node[loss] (task) at (112,80) {任务损失\\监督与回放};
  \node[loss] (align) at (112,56) {对齐损失\\统计或对应};
  \node[loss] (behavior) at (112,32) {行为损失\\概率匹配};
  \node[loss] (penalty) at (112,8) {模型参数惩罚\\重要性加权};
  \node[box,text width=24mm,fill=FigureOutput!10] (update) at (151,44)
    {组合有效目标\\同一次\\模型参数更新};
  \draw[flow] (data) -- (prediction);
  \draw[flow] (pairs) -- (features);
  \draw[constraint] (teacher) -- (behavior-pair);
  \draw[constraint] (reference) -- (difference);
  \draw[flow] (prediction) -- (task);
  \draw[flow] (features) -- (align);
  \draw[flow] (behavior-pair) -- (behavior);
  \draw[flow] (difference) -- (penalty);
  \draw[BookInk,line width=0.8pt] (132,8) -- (132,80);
  \draw[BookInk,line width=0.8pt] (task.east) -- (132,80);
  \draw[BookInk,line width=0.8pt] (align.east) -- (132,56);
  \draw[BookInk,line width=0.8pt] (behavior.east) -- (132,32);
  \draw[BookInk,line width=0.8pt] (penalty.east) -- (132,8);
  \draw[flow] (132,44) -- (update.west);
\end{tikzpicture}

\caption{同一候选参数接收的任务、对齐、行为与位移约束。实线表示计算流，虚线表示本轮固定参照的约束；旧记忆可以同时供给回放标签和历史行为。当前表示、输出和参数随学习变化，教师、参考参数及重要性不随动。并列损失是可选组合，不要求每轮全部使用。}
\label{fig:ket-update-objective}
\end{figure*}

\FloatBarrier
\subsubsection{更新范围与起点的确定}
\label{sec:ket-update-range}

同样的目标还可能因参数范围受限而无法实现。现有特征仍区分物种、只发生标签先验变化时，可检验仅校正偏置；表示足够而边界不同，可重学预测头；纹理已被编码器丢弃时，头再大也不能从现有特征恢复它。部分层、全参数或附加少量参数分别扩大不同可达解集合，也带来不同旧能力风险和成本。第~\ref{sec:ket-transfer-parameter-adaptation}节与第~\ref{sec:ket-transfer-added-parameters}节展开具体适应机制，这里确定可变参数和起点怎样约束维护。

起点决定哪些经验已经在参数中。当前版本带有最近更新，早期快照可能避开后续失稳，随机初始化则需要重新学习旧参数曾提供的内容。取起始参考$\bm\theta^\star$和允许索引集合$\mathcal J$，对$j\notin\mathcal J$固定$\theta_j=\theta_j^\star$，只限定可达参数，不保证可达函数均保留旧任务。尤其早期快照可能缺少已经生效的修正与撤回，恢复前必须核对，不能先服务再补要求。

单条事实请求也没有免除定位和验证。局部微调在候选层上降低编辑损失，并用保留样本或位移惩罚限制外溢；候选位置可以来自诊断，但损失下降只说明该位置能影响输出，不证明事实仅存于那里。学习编辑器则先训练从请求到参数增量的映射，将一部分逐请求求解成本转为准备成本。两者都还需要验证修改后的语义范围。

模型编辑网络MEND（model editor networks）利用梯度因子构造增量。对一条全连接层输入，梯度可写为误差信号$\bm u$与层输入$\bm v$的外积，序列或批量则累计多个外积。编辑器处理因子后组合修改量，避免直接预测整个大矩阵。准备阶段先按请求生成内层修改，再以改后请求损失和无关样本局部性损失训练编辑器；应用阶段固定编辑器，计算新请求梯度与增量。效果依赖编辑训练分布及基座访问条件，训练编辑器、取得梯度和应用增量的成本必须分别记录\scite{ket-mitchell2021-mend}。

秩一模型编辑ROME（rank-one model editing）先干预并恢复隐藏状态，寻找影响事实回忆的位置，再修改选定多层感知器（multilayer perceptron, MLP）中的线性映射。记$\bm W\in\mathbb R^{d_v\times d_k}$，输入激活$\bm k$称键，输出$\bm v=\bm W\bm k$称值；二者是内部向量，不是主体字符串与答案文本。主体、上下文和层共同形成非零编辑键$\bm k_\star$。先优化该处目标值$\bm v_\star$，使后续网络倾向新答案并满足保留目标，再求矩阵增量。设旧键二阶统计$\bm C=\mathbb E[\bm k\bm k^\top]$正定，以旧响应的平均平方变化为代价，在满足新键值硬约束的增量中选择最小者：
\begin{equation}
\min_{\Delta\bm W}\operatorname{tr}
(\Delta\bm W\bm C\Delta\bm W^\top),
\qquad
(\bm W+\Delta\bm W)\bm k_\star=\bm v_\star.
\label{eq:ket-update-rome-problem}
\end{equation}
迹项等于$\mathbb E\lVert\Delta\bm W\bm k\rVert_2^2$，故它直接比较旧键上的线性响应，而非所有网络输出。令残差$\bm r=\bm v_\star-\bm W\bm k_\star$，每行都需满足与$\bm k_\star$的内积等于相应残差。带正定二次罚项的约束最小化使行增量沿$\bm C^{-1}\bm k_\star$，归一化后得到
\begin{equation}
\Delta\bm W=
\frac{\bm r(\bm C^{-1}\bm k_\star)^\top}
{\bm k_\star^\top\bm C^{-1}\bm k_\star}.
\label{eq:ket-update-rome}
\end{equation}
右乘$\bm k_\star$时分母与键内积相消，因而实现指定残差，且外积的秩至多为一。非零键及正定$\bm C$使分母为正；奇异或病态统计还需声明正则化。这里是未中心化二阶矩，不是CORAL协方差。内部线性约束成立只完成写入步骤，仍须验证非线性网络中的改写、时间条件和依赖推理\scite{meng2022rome}。

大量事实的模型编辑MEMIT（mass-editing memory in a transformer）将请求分批写入多个关键MLP层。在某层，把$q$个键按列组成$\bm K\in\mathbb R^{d_k\times q}$，残差组成$\bm R\in\mathbb R^{d_v\times q}$。批量键未必能同时满足硬约束，因此以平方拟合误差和旧键响应代价折中：
\[
\min_{\Delta\bm W}
\lVert\Delta\bm W\bm K-\bm R\rVert_{\mathrm F}^2
+\operatorname{tr}(\Delta\bm W\bm C\Delta\bm W^\top)
\]
对增量求导，驻点满足$\Delta\bm W(\bm C+\bm K\bm K^\top)=\bm R\bm K^\top$，矩阵可逆时得到$\Delta\bm W=\bm R\bm K^\top(\bm C+\bm K\bm K^\top)^{-1}$。$\bm C$已含保留权重，相关或冲突键通过同一求解相互影响。各层写入后须重算下游键并分配剩余残差，不能复制同一增量。批量规模增加键、残差和线性求解成本；一次批量与逐条连续写入也有不同轨迹，均需复测早先有效请求\scite{ket-meng2022-memit}。

参数数量、位移大小和语义范围分别描述不同对象。一个共享标量若影响两个头，只改它也能改变两个任务；大量参数若作保持函数的重参数化，参数变化多却可不改响应。小位移也可能让边界附近样本换类。因而保留要求必须落到输入及响应上，不能由“只改少量参数”推断局部性；下一小节将计算这种区别。

持续预训练也需固定范围再比较。只改输出层保持编码器函数，却可能限制领域关系学习；全模型训练扩大可达函数，须增加通用与历史任务的约束证据。逐轮保存层范围、学习率、词元量及损失比例，才能区分数据增加、质量变化和范围扩大。底层不动不保证旧能力全保留，因为上层和输出竞争也参与旧预测。

类别加入后，旧类还要与新类竞争，而不只彼此比较。偏差校正BiC（bias correction）先用新数据、旧样本与蒸馏学习主分类器，再固定网络，用独立保留且旧新类适当平衡的校正材料，拟合新类分数$\tilde a_k=\alpha a_k+\beta$。这一步针对训练份额不均引起的偏差，旧类不作同样变换；校正集也占用标签，最终验收须另留材料。旧类相对分数不变不保证累积类别决策不变。第~\ref{sec:ket-transfer-output-semantics}节确定扩展目标，本节处理扩展后的竞争与维护\scite{ket-wu2019-bic}。

已有任务暂时正确，还不能说明以后能继续学。长期不活跃单元、大权重或优化状态变化可以作为调查线索，可塑性仍须用同预算后续学习测量。持续反向传播（continual backpropagation）跟踪单元效用及成熟程度，重置部分成熟低效单元的输入权重，将输出权重置零并清相关统计。零输出限制随机新特征的即时影响，成熟期防止刚生成就再次替换。连接不变，故这是参数重置而非扩容；低效单元仍可能保护稀有旧情形，原研究也不将该机制作为解决遗忘的保证，重置后须复测旧任务\scite{ket-dohare2024-plasticity}。

若希望旧经验提供后续可组合的参数方向，还可采用高效终身学习算法ELLA（efficient lifelong learning algorithm）。它写$\bm\theta_k=\bm L\bm s_k$，固定维数的共享基$\bm L\in\mathbb R^{d\times r}$保存方向，$\bm s_k\in\mathbb R^r$选择组合。新任务先给出独立参数估计$\hat{\bm\theta}_k$和半正定局部曲率权重$\bm D_k$，再以稀疏罚项鼓励少用方向，使部分系数为零，同时限制二次拟合误差：
\[
\min_{\bm s}
(\hat{\bm\theta}_k-\bm L\bm s)^\top
\bm D_k(\hat{\bm\theta}_k-\bm L\bm s)+\mu\lVert\bm s\rVert_1
\]
其中$\mu\geq0$控制$\ell_1$稀疏罚项，目标尺度已计入$\bm D_k$。系数取得后，固定各任务最近的系数，以累积二次统计及基正则化更新$\bm L$；旧任务即便系数不变，参数$\bm L\bm s_k$也随基变化。这样减少反复读取全部旧样本和重求所有系数的需要，却须保存参数、曲率或累积统计，且只是局部近似更新，非原联合目标的精确求解\scite{ket-ruvolo2013-ella}。

此处基维数和连接不增，改变的是共享参数。预测还需对应任务系数，身份未知时须另有合法选择规则。历史基若使新任务同效果所需样本减少，才支持前向帮助；基更新若使有效旧任务改善，才支持后向帮助。参数分解本身不提供这两项收益证据。

\subsubsection{更新计算的实施与控制}
\label{sec:ket-update-computation}

范围和起点确定后，才进入计算。阶段开始先保存固定教师、重要性、参数参照及有效请求；每步按声明权重取新旧批量，计算输出、表示和候选增量，只改允许参数。保留条件可要求投影实际方向或拒绝候选步。此时得到的是候选版本，仍须独立验收；训练目标下降只约束已使用材料，不等于任务或事实已修正。

令$\bm g^{(t)}$为当前目标梯度，$\eta>0$为步长，普通下降写作$\bm\theta^{(t+1)}=\bm\theta^{(t)}-\eta\bm g^{(t)}$。约束更新替换下降方向，直接编辑则可一次计算式~\eqref{eq:ket-update-rome}的增量。训练需重复前后向，投影另需旧梯度与子问题，编辑还可能需定位、统计或编辑器准备，成本应逐项记录。动量、二阶矩和随机状态帮助复现，也占存储，并可能依赖待撤回记录，不能只审参数文件。

用二维线性模型复算一次局部修改。教学上取$f_{\bm w}(\bm x)=\bm w^\top\bm x$、$\bm w=(1,1)^\top$，要求$\bm x_a=(1,0)^\top$从输出1改为2。允许两参数都变，以平方欧氏范数衡量改动，在满足请求的增量中选最小者：
\[
\min_{\Delta\bm w}\lVert\Delta\bm w\rVert_2^2,
\qquad
(\bm w+\Delta\bm w)^\top\bm x_a=2.
\]
约束固定$\Delta w_1=1$，取$\Delta w_2=0$最小化范数，故新权重为$(2,1)^\top$。在$\bm x_b=(0,1)^\top$上仍输出1，在$\bm x_c=(1,1)^\top$上却由2变3。这个增量满足目标且只改一个数，但并未保护所有非目标响应。

再保护$\bm x_c$输出2，就须加$\Delta w_1+\Delta w_2=0$，得增量$(1,-1)^\top$与新权重$(2,0)^\top$。平方代价由1升2，范数由1升$\sqrt2$，但$\bm x_b$变为0。三个要求联立不相容；同一$\bm x_a$同时要求1与2则更直接无解。增加优化步数不扩大可行集合。算例解释响应约束的关系，ROME实际使用的定位、内部键值与二阶统计仍须另行处理。

复杂旧任务通常不能写成几个线性输出条件，可以改用记忆损失的局部变化。记旧任务损失$L_k$、梯度$\bm g_k=\nabla L_k(\bm\theta)$，候选方向$\bm v$对应位移$-\eta\bm v$。局部二阶余项有界时，Taylor展开给出
\begin{equation}
L_k(\bm\theta-\eta\bm v)
=L_k(\bm\theta)-\eta\bm g_k^\top\bm v
+O(\eta^2\lVert\bm v\rVert_2^2).
\label{eq:ket-update-gradient-taylor}
\end{equation}
一阶损失变化为$-\eta\bm g_k^\top\bm v$，所以保护条件是内积非负。梯度情景记忆GEM（gradient episodic memory）对每个旧任务计算该条件，在所有可行方向中选择离当前梯度$\bm g$最近的一项：
\begin{equation}
\min_{\bm v}\frac12\lVert\bm v-\bm g\rVert_2^2,
\qquad \bm G^\top\bm v\ge\bm0,
\quad
\bm G=(\bm g_1,\ldots,\bm g_K).
\label{eq:ket-update-gem}
\end{equation}
$\bm G$按列保存$K$个旧任务梯度，平方距离尽量保留当前方向，线性约束排除一阶损害记忆的分量。零方向总可行，原梯度若已可行就无须投影，否则解凸二次规划。保护对象只是记忆经验损失，不是未知总体风险\scite{ket-lopezpaz2017-gem}。

参数维数远大于$K$时，直接求方向变量成本较高。对每个不等式引入非负乘子$\bm\alpha\in\mathbb R^K$，写Lagrange函数：
\[
\mathcal L(\bm v,\bm\alpha)
=\tfrac12\lVert\bm v-\bm g\rVert_2^2
-\bm\alpha^\top\bm G^\top\bm v.
\]
关于$\bm v$的驻点给出$\bm v=\bm g+\bm G\bm\alpha$；消去参数维变量，并将对偶最大化改写成最小化，得到
\begin{equation}
\min_{\bm\alpha\ge0}
\tfrac12\bm\alpha^\top\bm G^\top\bm G\bm\alpha
+\bm\alpha^\top\bm G^\top\bm g,
\qquad
\bm v^\star=\bm g+\bm G\bm\alpha^\star.
\label{eq:ket-update-gem-dual}
\end{equation}
该式只优化$K$个乘子，再恢复方向。例如教学上取$\bm g=(-1,1)^\top$、$\bm g_1=(1,0)^\top$，要求$v_1\geq0$，故投影为$(0,1)^\top$，$\alpha^\star=1$。去掉损害旧任务的一阶分量，也改变了当前下降方向，不能期待当前学习速度完全不变。

投影只控制一阶项。若$L_k$的梯度满足$\beta_k$-Lipschitz连续，下降引理给出有限步长的上界：
\[
L_k(\bm\theta-\eta\bm v)
\le L_k(\bm\theta)-\eta\bm g_k^\top\bm v
+\frac{\beta_k\eta^2}{2}\lVert\bm v\rVert_2^2.
\]
由上界保证不增，需要$\bm g_k^\top\bm v\geq\beta_k\eta\lVert\bm v\rVert_2^2/2$，比一阶非负更强；边界内积为零时损失仍可能上升。记忆估计与求解误差进一步限制保证，可减步长、复算损失或拒绝候选，却不能据投影宣称全程无遗忘。动量和预条件还会改变位移，因此真正受约束的应是优化器执行方向，而不只是原梯度。

元经验回放MER（meta-experience replay）利用同一局部关系鼓励经验间相容。先沿经验$i$下降，会使$j$损失近似变为$L_j(\bm\theta)-\eta\bm g_j^\top\bm g_i$，正内积支持局部共同下降，负值提示干扰。MER混合当前与回放经验做内部更新，再在批内、批间向学后参数插值，即采用Reptile式元更新近似促进这种关系。它没有逐任务执行GEM的硬投影，局部内积改善也不证明未见任务更容易；前、后向收益仍须对照\scite{ket-riemer2018-mer}。

多项约束可能把方向压近零，使新内容难以学入。应先判断要求是否过时或允许折中，不能事后删除受损任务来提高平均分。放松约束针对保留冲突，扩大参数范围针对可达解不足，补证据针对识别缺口；重置优化状态只改变后续动力学，不直接补回参数已丢的信息。失败诊断决定采用哪项操作。

若使用协议允许当前测试输入参与学习，还可在预测过程中更新。\term{测试时适应}（test-time adaptation, TTA）利用这些输入改变协议允许的参数或统计状态，使后续预测依赖更新后的模型。Tent以测试熵最小化（test entropy minimization）为目标：在当前批量$B$上计算类别概率的不确定性，并减小平均预测熵：
\[
L_{\mathrm{ent}}(\bm\theta;B)
=-\frac1{|B|}\sum_{\bm x\in B}\sum_k
p_{\bm\theta}(k\mid\bm x)\log p_{\bm\theta}(k\mid\bm x)
\]
其中$p_{\bm\theta}(k\mid\bm x)$为当前类别概率，$B$非空；降低熵让概率更集中，不直接降低真实错误。典型实现只训练归一化层仿射参数，并以当前批量计算统计，其余权重冻结。原始在线计算先产生本批输出再更新，影响下一批；若重新回答本批，须增加前向且另声明计分协议。标签不可见不等于输入未用于学习\scite{wang2021tent}。

取无标签流$B^{(1)},B^{(2)},\ldots$，先预测后更新的状态关系为
\[
\bm\theta^{(t)}
\ \longrightarrow\ p^{(t)}(B^{(t)})
\ \longrightarrow\ L_{\mathrm{ent}}^{(t)}
\ \longrightarrow\ \bm\theta^{(t+1)}
\ \longrightarrow\ p^{(t+1)}(B^{(t+1)}).
\]
若阴影乙已被误判为甲，熵下降可能继续推高甲概率，下一批再从偏移后的模型形成目标。恒预测甲的熵可以很低，却对乙始终错误。单样本统计较弱，批量统计利用更多输入，却带来等待与样本间依赖；持续状态既能累积适应，也能累积错误。按环境重置须有许可边界，不能看结果后才划环境。

锐度感知的可靠熵最小化SAR（sharpness-aware and reliable entropy minimization）针对小批量、混合域和类比例变化筛选可靠输入，并考察参数邻域中的熵损失，减少对敏感局部方向的依赖。它还监测运行熵统计，在塌缩迹象下恢复状态；单类塌缩可能低熵，故不是简单“高熵重置”。阈值、更新范围与恢复规则预先固定，准确性仍由真实反馈检验。恢复权重改变函数，清优化器改变以后如何学，两项要分别记录\scite{ket-niu2023-sar}。

不能区分输入故障与真实规律变化时，继续降低代理损失可能扩大偏差。此时停止适应、使用仍合规的最近验收版本或请求标签都是候选，不应把更新本身当作目标。整批适应、严格在线与可回看的离线协议须分别报告；无标签权限不授权未来输入和任意重排。

需要广泛重估或处理撤回时，可以重新执行训练。令$\mathcal A(S;\xi)$为指定学习过程，$\xi$包含初始化、顺序及随机操作，按记录身份取撤回集合$U$和$S^-=S\setminus U$。参照$\mathcal A(S^-;\xi')$应重建受影响的预处理、特征和状态，不能保留依赖$U$的未处理缓存。普通适应也可重训，但只有排除相关影响路径并满足规定参照，才支持机器遗忘。

教学上设记录$u$第一次进入训练块$S_2$，$C_j$保存处理完块$j$的完整状态；对照原轨迹与排除它的重训轨迹：
\[
\begin{array}{rcl}
\text{原轨迹}&:&
C_0\xrightarrow{S_1}C_1
\xrightarrow{S_2\ni u}C_2
\xrightarrow{S_3}C_3,\\
\text{完整重训}&:&
C'_0\xrightarrow{S_1}C'_1
\xrightarrow{S_2\setminus\{u\}}C'_2
\xrightarrow{S_3}C'_3.
\end{array}
\]
从头重训重新计算全部块；只有$C_1$连同预处理与随机状态都不依赖$u$，才可考虑复用前缀。由$C_2$继续训练仍带着既有影响。删记录还可能改变随机操作对应关系，相同种子不保证逐步对齐；验收须说明比较一次结果还是算法的输出分布。

SISA将训练组织为分片、隔离、切片与聚合（sharded, isolated, sliced, aggregated），以事前隔离减少撤回重算。互斥分片各自训练，片内切片顺序处理并留快照；若$u$在分片$a$第二切片，从此前状态重算该片后缀，其他片不动：
\[
\begin{array}{rcl}
\text{分片 }a&:&
C_{a,0}\to C_{a,1}
\xrightarrow{S_{a,2}\ni u}C_{a,2}\to C_{a,3},\\
\text{撤回后 }a&:&
C_{a,1}\xrightarrow{S_{a,2}\setminus\{u\}}C^-_{a,2}
\xrightarrow{S_{a,3}}C^-_{a,3},\\
\text{分片 }b&:&C_{b,3}\quad\text{保持，随后重新聚合。}
\end{array}
\]
随后重新聚合。少算的部分来自隔离依赖，代价是子模型、快照和聚合资源，任务效果也可能不同。参照须是同一分片、切片及聚合规则排除$U$的学习过程，不是全部剩余数据上的单体训练。共享预处理若看过所有数据，仍可能留下跨片影响，须纳入重建\scite{bourtoule2021sisa}。

另一条路线以现有参数近似删除后的解。认证数据删除（certified data removal）在光滑强凸正则线性模型等特定条件下，用剩余目标$L^-$的梯度和Hessian作Newton修正，再将残差界与预先加入的随机扰动相连。一步Newton更新写为$\bm\theta^-=\bm\theta-[\nabla^2L^-(\bm\theta)]^{-1}\nabla L^-(\bm\theta)$，须保证曲率可逆并控制近似误差。小残差只支持优化接近，不独立提供发布分布保证；任意神经网络的反向更新不能沿用认证。完整条件见\BookSectionRef{sec:privacy-unlearning-methods}{卷3《基础、理论与可信性》第16章“隐私保护”第16.4.2节“训练影响的消除方案”}\scite{guo2020certifiedremoval}。

把撤回照片分错，可能只提高该照片损失，中间特征仍保存其背景，对其他预测的影响也可能仍在；拒答规则更只约束接口。因此放弃某项能力、抑制答案和消除训练影响是不同目标。重建后还须继承其他有效编辑与行为要求，再按第~\ref{sec:ket-update-unlearning-validation}节检查撤回对象、剩余任务和重训练参照。

\subsection{模型结构的调整}
\label{sec:ket-update-structure}

若当前结构没有同时满足新旧要求的函数，延长训练不会扩大可达集合；共用参数中的干扰也可能让两项要求难以兼顾。此时可增减组件、改变连接及允许的调用路径，让参数学习在另一种配置中进行。固定连接下的权重重置仍是参数操作，新增模块的权重也仍需前节的证据和目标。结构改变提供候选容量，并不直接证明已获得能力。

\subsubsection{组件配置与容量调整}
\label{sec:ket-update-capacity}

为隔离新梯度，可以为各任务保留独立网络；旧函数不受新网络训练影响，但参数随任务数增长，新经验也未直接反馈旧网络。另一种配置共享部分层并保留专属分支，以减少重复，代价是共享处仍可干扰。共享底座若冻结，限制新任务可读取的信息；若更新，多条旧路径都可能变。比较这些配置须固定任务和预测接口，尤其不能免费提供某方法才需要的任务身份。

渐进网络（progressive networks）在隔离新参数的同时，允许新列读取旧列特征。它为任务$k$增加一列，冻结已有列，用横向连接引入旧表示。记第$\ell$层表示为$\bm h_k^{[\ell]}$，忽略偏置的概括形式为
\[
\bm h_k^{[\ell]}
=\sigma\!\left(
\bm W_k^{[\ell]}\bm h_k^{[\ell-1]}
+\sum_{j<k}\bm U_{k,j}^{[\ell]}\bm h_j^{[\ell-1]}
\right).
\]
这里$[\ell]$标层而非阶段，$\bm W_k^{[\ell]}$为新列层间矩阵，$\bm U_{k,j}^{[\ell]}$将旧列$j$的特征映到新层，$\sigma$是非线性激活。只学习新列与横向权重；旧任务沿不变旧列预测，新任务还需计算被读取的旧列。因而保留依赖旧调用路径也不变，新知识并未回写旧列。列、连接和推理计算可随任务增长，成本不能只计最新列\scite{ket-rusu2016-progressive}。

扩容也可服务原任务，例如RGB与深度输入使用不同编码模块，但仍识别同一物种集合。输出或成功标准若新增，才先交第~\ref{chap:ket-transfer-adaptation}章明确目标。冻结组件只保证其参数数值不变，输入处理、统计、共享层、路由或校准变化仍能改变端到端结果，须分别验证。

例如，温室分支保持正确，路由却在新增雨季分支后把所有绿背景输入送去雨季路径。温室也有绿背景照片时，错误分配便可能降低温室准确率，旧分支参数完全不动也无济于事。这需要分别检查组件函数与实际调用；节点增加只说明容量改变，不证明正确能力能够被请求使用。

分支持续增长后，还需决定何时整合或回收。若多个分支在同样输入和误差口径下实现近似功能，可调查整合；旧传感器退役时，可调查对应模块是否仍被其他任务依赖。扩容继续提供隔离，重分配改变有限容量的服务对象，压缩减少重复却引入逼近误差。第~\ref{sec:ket-integration-composition}节和第~\ref{sec:ket-integration-distillation}节说明组合与蒸馏机制，本章确定触发证据、保留接口和验收条件。

回收先追踪任务、条件、横向连接和版本依赖。稀少调用可能对应必须服务的罕见物种，不能直接当作无价值。合并后须检查输入输出、类别顺序和尺度兼容，更新旧路由引用，再测试后续学习。存储减少也要减去替代模型、转换器及必需快照的新增占用。

\subsubsection{组件连接与路由安排}
\label{sec:ket-update-routing}

温室和雨季组件都有以后，一次输入仍须选到适用路径。\term{路由}（routing）就是用当前可见信息决定调用哪些组件的规则。可信相机标识可选输入模块，环境未知时则需从输入或上下文估计。任务身份说明要完成什么，环境身份说明条件，类别却是待推断答案；向路由器提供真实类别或协议不可见的任务号，会改变问题，而不只是改进路由。

对已保留的$f_1,f_2$，硬路由$r(\bm x)\in\{1,2\}$只取一个输出。若两者对同一类别集合给概率，也可由输入相关权重混合：
\[
p(y\mid\bm x)=
\alpha_1(\bm x)p_1(y\mid\bm x)
+\alpha_2(\bm x)p_2(y\mid\bm x),\qquad
\alpha_j\ge0,\quad\alpha_1+\alpha_2=1.
\]
$p_j$是组件概率，非负且和为一的$\alpha_j$保持混合结果为概率。硬选择只有在不运行未选路径时才省对应计算，错选则绕过适用能力；混合需多组件前向，还须核对类别含义与校准。不同温度logit不能因同维就直接相加。共享底座变化会同时改变组件及路由的输入，兼容性也须复查。

学习提示L2P（learning to prompt）用提示池在冻结预训练模型上作输入相关选择。提示是参与模型计算的可学习附加向量，每条配一个检索键；固定查询函数先编码输入，再按相似度选择提示并参与预测。任务损失学习提示及相关预测参数，键匹配损失学习选择关系。推理不必读取真实任务身份，但查询前向、带提示前向、键和池容量仍有成本。错选或共享提示被覆盖，也可能破坏旧能力\scite{ket-wang2021-l2p}。

路由权重的学习仍属于参数更新，结构安排则规定可读信息、可达组件及是否允许并行调用。既有图因输入选择不同路径，不等于每次都重设计结构。图~\ref{fig:ket-update-operations}从同一双组件模型比较改权重、增节点连接和选一次路径，可据此分别核算成本与风险。

\begin{figure*}[htbp]
\centering
% The shared initial graph has two components A and B.
% Added nodes/edges encode capacity; line emphasis encodes one invocation.
\begin{tikzpicture}[
  x=1mm,y=1mm,font=\figurefont\small,
  box/.style={rectangle,draw=BookInk,fill=BookPaper,line width=0.7pt,
    align=center,inner sep=1.5mm,minimum height=8mm},
  module/.style={box,fill=FigureModel!10,minimum width=12mm},
  flow/.style={draw=BookInk,line width=0.7pt,-{Latex[length=1.7mm]}},
  unused/.style={flow,dashed,draw=BookMuted},
  active/.style={flow,line width=1.2pt,draw=BookTeal}
]
  \node[anchor=west,text=BookInk] at (0,122) {共同初始模型：已有组件A、B及两条可选路径};
  \node[box] (start) at (13,105) {预测输入};
  \node[box] (router) at (52,105) {路由};
  \node[module] (a) at (103,114) {A};
  \node[module] (b) at (103,96) {B};
  \node[box] (out) at (150,105) {预测输出};
  \draw[flow] (start) -- (router);
  \draw[flow] (router) -- (a);
  \draw[flow] (router) -- (b);
  \draw[flow] (a) -- (out);
  \draw[flow] (b) -- (out);
  \draw[BookRule,line width=0.6pt] (0,85) -- (163,85);

  \begin{scope}
    \node[anchor=west,font=\figurefont\small\bfseries] at (0,77)
      {A\quad 模型参数变化};
    \node[box] (pi) at (25,65) {预测输入};
    \node[box] (pr) at (25,50) {路由};
    \node[module,fill=FigureOutput!15] (pa) at (12,31) {A$^\star$};
    \node[module] (pb) at (38,31) {B};
    \node[box] (po) at (25,12) {预测输出};
    \draw[flow] (pi) -- (pr);
    \draw[flow] (pr) -- (pa);
    \draw[flow] (pr) -- (pb);
    \draw[flow] (pa) -- (po);
    \draw[flow] (pb) -- (po);
    \node[text=BookMuted,font=\figurefont\footnotesize] at (25,0)
      {仅A的权重改变};
  \end{scope}
  \begin{scope}[xshift=56mm]
    \node[anchor=west,font=\figurefont\small\bfseries] at (0,77)
      {B\quad 结构变化};
    \node[box] (si) at (25,65) {预测输入};
    \node[box] (sr) at (25,50) {路由};
    \node[box,fill=FigureModel!10] (sa) at (8,31) {A};
    \node[box,fill=FigureOutput!15] (sc) at (25,31) {C};
    \node[box,fill=FigureModel!10] (sb) at (42,31) {B};
    \node[box] (so) at (25,12) {预测输出};
    \draw[flow] (si) -- (sr);
    \draw[flow] (sr) -- (sa);
    \draw[flow] (sr) -- (sb);
    \draw[active] (sr) -- (sc);
    \draw[flow] (sa) -- (so);
    \draw[flow] (sb) -- (so);
    \draw[active] (sc) -- (so);
    \node[text=BookMuted,font=\figurefont\footnotesize] at (25,0)
      {新增C及其连接};
  \end{scope}
  \begin{scope}[xshift=112mm]
    \node[anchor=west,font=\figurefont\small\bfseries] at (0,77)
      {C\quad 调用变化};
    \node[box] (ri) at (25,65) {预测输入};
    \node[box] (rr) at (25,50) {路由};
    \node[module] (ra) at (12,31) {A};
    \node[module] (rb) at (38,31) {B};
    \node[box] (ro) at (25,12) {预测输出};
    \draw[active] (ri) -- (rr);
    \draw[active] (rr) -- (ra);
    \draw[unused] (rr) -- (rb);
    \draw[active] (ra) -- (ro);
    \draw[unused] (rb) -- (ro);
    \node[text=BookMuted,font=\figurefont\footnotesize] at (25,0)
      {本次选A，B未调用};
  \end{scope}
\end{tikzpicture}

\caption{同一初始模型上的三种变化。左侧只改组件权重，中间新增组件C及其可达连接，右侧在既有连接中以粗实线选中A、虚线标未调用B。结构容量由节点和连接增减表达，框面积不代表参数量。操作可以组合，路由规则的权重也可学习。}
\label{fig:ket-update-operations}
\end{figure*}

\subsection{外部知识的更新}
\label{sec:ket-update-external}

负责人任命带来源和时间，外部记录能够直接保存这些字段，因而可与参数编辑并列作为维护候选。但记录变了，回答未必变：检索可能还读旧缓存，生成器也可能忽略新证据。外部更新须同时完成内容与有效范围修改，以及实际调用同步。

\subsubsection{记录内容与有效范围的修改}
\label{sec:ket-update-records}

检索增强生成RAG（retrieval-augmented generation）先按问题取得外部文档，再让生成器在这些材料条件下回答。参数承担语言与任务计算，外部库提供可替换内容；更换文档或索引可改变生成条件，而不必每次重训生成器。参数旧事实仍可能与检索证据竞争，生成也未必忠实于材料。该分工提供更新接口，不自动处理有效时间和撤回\scite{ket-lewis2020-rag}。

若希望不改底模而保存局部编辑，可以采用SERAC（semi-parametric editing with a retrieval-augmented counterfactual model），即用检索增强的替代模型承接编辑。缓存保存请求，范围分类器判断当前输入是否受影响：范围外调用底模，范围内按相关请求生成修正回答。准备阶段学习范围与修正生成，新增请求时写缓存，不直接改底模。局部性转而依赖选择是否准确，漏选保留旧答，过宽选择误改无关输入，缓存增长又增加检索与冲突成本\scite{ket-mitchell2022-serac}。

甲所记录可按实体、关系、对象、有效区间、来源和状态保存。新增“乙自7月1日起负责”，还须把甲的有效终点收束到同日；否则两记录仍可能同时适用。直接替换可维护仅含当前状态的副本，却丢历史依据；保留两段时间才能分别答6月和7月。过期表示不适用于当前时间，不代表历史为假；物理删除则使副本不可访问，作用不同。

同名主体必须按标识区分，时间也有发布、生效和入库三种角色。7月3日入库的文件可能7月1日生效，较晚入库的旧报道不能因此覆盖任命。来源相矛盾而优先级无依据时，应保存条件与分歧并让结论待决；按入库顺序只留最后一条，不能消除真实冲突。

删除还要追踪原文生成的摘要、向量、缓存与索引。主记录不存在，派生副本仍可能供给同一内容。合法历史事实仅不再当前适用时，可保留历史区间；请求禁止使用特定记录时，则按要求清理或隔离副本。若它曾进入训练，还须另行处理参数影响，外部删除不能隐含完成机器遗忘。

\subsubsection{调用结果与生成行为的同步}
\label{sec:ket-update-sync}

回答依次经过问题解析、候选检索、范围筛选和生成，每个环节都可能使用旧版本。解析应取得主体与查询时间，候选可包含两代事实，再按区间及来源选适用证据，生成器据此回答。文本最相似的条目不必时间适用；“现任”还须绑定服务时的当前日期，而非模型训练日期。

任命生效后，先关联新旧记录并修改有制度依赖的当前签批和办公室摘要，再使索引可检索、缓存失效或绑定正确时间。7月15日排除甲的任职记录，6月15日仍选它；历史审批签字人不属于“当前负责人”的派生事实，不能随新任命重写。内容修改与调用同步由此具有不同完成条件。

缓存只以问题文字为键时，任命前后的“谁负责”可能混用旧结果。可按记录依赖失效，或绑定记录版本和时间条件，具体实现不同但都须验证适用性。模型、记录与索引未同步的中间状态，应暂停相关服务或明确证据状态，不能算完整验收。索引及缓存工程留给应用与\BookPartRef{part:systems}{卷12《系统》第1部分}，这里规定需保持的语义关系。

两条证据冲突未决时，回答也应保留条件和不确定性，不能按新旧次序抹去分歧。范围模型不识别历史日期、生成器不使用证据，则应补相应训练；反复添同一事实不能直接修正这个环节。记录候选、选中证据和最终答案，才能把失败定位到检索、范围或生成。

检索为空时，底模仍可能从参数答出已删内容，所以外部不可检索不是训练影响消除证据。参数重训也不自动清理缓存。请求涉及两条路径时，分别执行并验证，再将相容的模型、记录、索引和缓存组成系统候选。

\section{知识更新结果的验证}
\label{sec:ket-update-validation}

候选版本已经生成，现在按第~\ref{sec:ket-update-goals}节的要求分别验收：原任务在当前域是否改善，编辑是否在正确条件及依赖中生效，旧能力和后续学习是否保持，撤回是否满足指定参照。四项主张观察不同对象，同一版本可能只通过其中几项；汇总分数不能抵消未满足的硬要求。

\subsection{当前任务表现的验证}
\label{sec:ket-update-current-validation}

判断雨季预测是否改善，先固定任务、数据使用方式、访问权限与调参机会，报告实际训练和推理预算。原版本提供更新前结果；不更新基线则将原权重放进相同解码、接口和评测流程，区分外围修复与模型更新的贡献。选过学习率的数据不能再称独立最终验收集，划分协议见第~\ref{sec:ket-evaluation-splits}节。

雨季与仍服务的温室应分别计分，再按物种、遮挡和相机检查。教学上取两域原风险$(0.1,0.4)$、新风险$(0.2,0.2)$，风险越低越好。等权平均从$0.25$降至$0.2$，却伴随温室风险翻倍；若要求温室不退化，就不能接受。允许折中时也须事先声明域权重，不能只展示有利平均值。

深度模型的验收还要检查颜色区分类、缺测和配准失败，而不只看总体特征误差。它们可以暴露输入信息缺失与对应错误，但机制判断还需相应样本证据。回归按单位和条件误差检查，时序按可见窗口与预测时距分组，图数据按节点邻域依赖分析；高度相关记录不能充作独立重复证据。

无目标标签时，协方差距离、域预测准确率、熵和增广一致性只能检验代理。域准确率的随机基准还取决于评价采样比例，不能省略该条件。没有真值或等价任务反馈，就应报告代理变化和假设检查，同时明确真实风险验收未完成。少量标签结果也需结合重复植株、时间相关及运行变化报告不确定性，不能由一次划分推广全部未来域。

\subsection{修改范围与关联一致性的验证}
\label{sec:ket-update-edit-validation}

编辑测试从请求及其规则生成，不从已答对的问题反选。指定表达问“2026年7月15日甲所负责人是谁”，预期乙；等义改写可以问“甲所在2026年7月15日由哪位负责人任职”，其他语言表达另行核对语义。指定表达失败表示目标未实现，只有原提示成功则提示表达泛化不足。近重复提示数量不能替代新的表达覆盖。

依赖组按给定制度问7月15日签批人和联系办公室，答案应随任命指向乙；保留组检查成立年份、同名他所、6月15日负责人和历史审批签字人。直接事实成功而依赖组失败，说明相关使用仍未修正；保留组失败，则需检查身份、时间或无关外溢。局部性按请求及依赖定义，不要求所有相邻语义都不变。

多跳编辑问答基准MQuAKE（multi-hop question answering for knowledge editing）将多个事实组合成问题，检验直接事实成功之后能否得到新的链式答案\scite{ket-zhong2023-mquake}。RippleEdits则按逻辑泛化、组合、实体别名及应保留关系区分编辑传播；其“传播”指哪些依赖输出应随编辑变，而不是任意关联问题都该变\scite{ket-cohen2024-ripple}。

这两类测试用于甲所时，仍须先声明制度依赖并制定答案。跨语言测试检查翻译保真，时间测试固定查询日期，不能把不同权限、不同历史条件的问法算成同一事实改写。基准提供测试维度，不替代当前请求的有效范围。

批量数、序列长度和依赖复杂度分别控制不同干扰。十条独立请求、一条负责人连续改十次和十关系推理链，不应合为同一压力条件。可固定每批数增加阶段，或固定阶段数改变依赖范围，逐轮复查仍有效请求。后续合法任命替代当前答案时，按新条件计分，不把必要覆盖算遗忘，也不只测最后一次编辑。

外部更新按内容、调用和响应定位失败。写入成功却命中旧缓存，是同步问题；候选缺新记录，需查索引或召回；已取得两代记录仍按入库日期回答历史，则需查范围或生成。保存模型、库、索引、缓存版本及选中证据，才能判断第~\ref{sec:ket-update-sync}节的操作是否完成。写入接口成功不能替代端到端验收。

\subsection{能力保留与后续学习的验证}
\label{sec:ket-update-retention-validation}

连续吸收可能在中途破坏旧能力，故需逐阶段复测。按任务进入顺序取$\mathcal T_k$，令$a_k^{(t)}$为完成阶段$t$后，在固定评价协议及相容有效要求下测得的准确率等指标，越大越好。记录可以排列为
\[
\begin{array}{c|ccc}
 & k=1 & k=2 & k=3\\ \hline
t=1&a_1^{(1)}&\text{--}&\text{--}\\
t=2&a_1^{(2)}&a_2^{(2)}&\text{--}\\
t=3&a_1^{(3)}&a_2^{(3)}&a_3^{(3)}
\end{array}
\]
同一行比较当前版本的各任务，同一列追踪任务跨阶段变化。未测标记不是零；协议允许学前零样本评价时，才填实际观测。每格需附样本、预算和计分时点，通用矩阵见第~\ref{sec:ket-evaluation-continuous}节。

保留可参照学完任务的$a_k^{(k)}$，也可参照此前最佳值，二者回答不同下降幅度，须声明基准，见第~\ref{sec:ket-evaluation-backward}节。负责人任命变化后，“现任”旧答案甲不能继续作真值，应将查询时间显式化；撤回后禁止维持的行为也应移出保护目标。只有任务、时间和请求含义相容，跨阶段下降才能解释为对应退化。

测可塑性则需继续训练，而非只测旧任务。原、更新版本各复制独立起点，在同样新内容上使用相同样本顺序、标签、可变层和更新次数，记录各预算的损失。相同旧分数可能对应不同新学习曲线。若比较达到风险阈值$\delta$的样本数$N_\delta$，应预先固定阈值和实测预算；未过线只记预算内未达标，不外推到达量。

知识积累另需控制相关历史经验。容量、当前数据、预算及其他预训练因素可比时，有该经验的起点若同样样本下风险更低，或同阈值所需样本更少，才支持前向帮助；它与学前零样本前向迁移不同，见第~\ref{sec:ket-evaluation-forward}节。后向帮助则在不额外增加旧监督的条件下检查有效旧任务是否改善。多保存模型或扩大容量都不是经验贡献的独立证据。

表示更新还须区分探测与使用。后来用标签拟合探测头，测得的是读出阶段可利用的信息；原有任务头当时的输出，才对应即时任务能力。探测头的数据或计算不一致，还会引入新混杂。图~\ref{fig:ket-update-evaluation}将旧复测和后续训练分支隔离，分别观察保留与可塑性；若要判断积累，还需增加历史经验有无的对照。

\begin{figure*}[htbp]
\centering
% Protocol diagram only: no invented accuracy or learning-curve data.
% Each evaluation uses a separate copy; new learning cannot alter old retests.
\begin{tikzpicture}[
  x=1mm,y=1mm,font=\figurefont\small,
  box/.style={rectangle,draw=BookInk,fill=BookPaper,line width=0.7pt,
    align=center,inner xsep=2mm,inner ysep=2mm,minimum height=12mm},
  flow/.style={draw=BookInk,line width=0.8pt,-{Latex[length=2mm]},
    rounded corners=1mm}
]
  \node[box,text width=60mm,fill=FigureModel!10] (original) at (40,93)
    {原版本\\冻结评测起点，复制两份};
  \node[box,text width=60mm,fill=FigureOutput!10] (updated) at (124,93)
    {更新版本\\冻结评测起点，复制两份};
  \node[box,text width=29mm] (old-a) at (20,62)
    {旧任务复测\\当前有效要求};
  \node[box,text width=29mm] (new-a) at (60,62)
    {同一新任务学习\\等样本、等预算};
  \node[box,text width=29mm] (old-b) at (104,62)
    {旧任务复测\\当前有效要求};
  \node[box,text width=29mm] (new-b) at (144,62)
    {同一新任务学习\\等样本、等预算};
  \draw[flow] (original.south) -- (40,79) -| (old-a.north);
  \draw[flow] (40,79) -| (new-a.north);
  \draw[flow] (updated.south) -- (124,79) -| (old-b.north);
  \draw[flow] (124,79) -| (new-b.north);
  \node[box,text width=29mm,fill=FigureInput!10] (old-ra) at (20,33)
    {原版本的\\旧任务结果};
  \node[box,text width=29mm,fill=FigureLoss!10] (new-ra) at (60,33)
    {原版本的\\后续学习过程};
  \node[box,text width=29mm,fill=FigureInput!10] (old-rb) at (104,33)
    {更新版本的\\旧任务结果};
  \node[box,text width=29mm,fill=FigureLoss!10] (new-rb) at (144,33)
    {更新版本的\\后续学习过程};
  \draw[flow] (old-a) -- (old-ra);
  \draw[flow] (new-a) -- (new-ra);
  \draw[flow] (old-b) -- (old-rb);
  \draw[flow] (new-b) -- (new-rb);
  \node[align=center,text=BookInk,text width=160mm] at (82,6)
    {两侧旧任务结果对照：能力保留；两侧同预算学习对照：可塑性\\
     进一步控制相关历史经验的有无：检验跨阶段学习收益};
\end{tikzpicture}

\caption{旧能力与后续学习的分离对照。两个版本各复制为旧任务复测与新任务学习分支；旧复测使用相容有效要求，新学习使用相同样本顺序和预算，不能用学后状态覆盖学前复测。历史经验的贡献另需有无对照。图展示协议，不含实测曲线或假设分数。}
\label{fig:ket-update-evaluation}
\end{figure*}

\FloatBarrier
\subsection{训练影响消除的验证}
\label{sec:ket-update-unlearning-validation}

删除验收先固定记录集合$U$、剩余集$S^-$、学习过程$\mathcal A$和观察者能取得的对象。任务、预处理、模型族及训练协议相同，才可生成对应重训练参照；SISA采用同一分片过程。候选和参照还须以相同规则落实有效事实及行为请求，否则差异可能来自版本要求。随机状态和派生物的范围也属于定义，不只是被删除文件名。

撤回样本损失、剩余任务风险和参照距离分别提供证据。损失升高可能只来自输出抑制，剩余任务正确也只表示未观察到该类损害。给定输入集上的预测差可暴露残留，但不约束集外；随机重训还需估计自然运行差异，不能把两次随机差全部归于删除失败。

参数距离还依赖参数化，隐藏单元置换可改变参数而不改函数，小参数差也能改敏感输入。输出距离只约束观察输入。若主张算法发布结果的分布保证，就需固定$Z_{\mathrm{del}}$和$Z_{\mathrm{retrain}}$所含模型及可见状态，并比较所有可测事件$E$，而不只平均几个测试输出。例如，一种近似保证要求
\[
\mathbb P(Z_{\mathrm{del}}\in E)
\le e^\varepsilon\mathbb P(Z_{\mathrm{retrain}}\in E)+\delta,
\]
以及交换两侧后的不等式。$\varepsilon\geq0,\delta\geq0$限制可区分程度，两侧概率来自训练及删除随机性；参数须由具体机制保证，而非少量提示一致率估出。若观察者还能看旧模型或日志，就需比较联合观察，不能用新模型边缘分布撤销既有披露。不同方法的误差定义须保留，完整边界见\BookSectionRef{sec:privacy-unlearning-validation}{卷3《基础、理论与可信性》第16章“隐私保护”第16.4.4节“撤回结果与效率的验收”}。

成员推断测试是否能识别某记录的训练参与，其他攻击也可暴露特定残留；失败只说明既定能力、样本和预算未检出。机制若要求强凸、扰动或保存状态，缺少前提就不能沿用认证。经验结果、条件核查和证明应分列，拒答或清缓存不构成发布分布保证。

证据满足要求时接受候选，否则继续修正或不启用。硬保证缺失不能由其他任务高分抵消；剩余能力受损时，返回范围、参照或算法检查。旧状态若含已撤回影响，也不能作为不受约束的回退。表~\ref{tab:ket-update-validation}把各主张与证据并列，用于检查指标是否回答了当前问题。

\begin{table*}[htbp]
\small
\begin{tabularx}{\linewidth}{@{}p{19mm}XXX@{}}
\toprule
待验证主张 & 比较基线 & 主要观测对象 & 常见误判 \\
\midrule
当前效果改善 & 原版本与同协议不更新基线 & 当前及仍服务原环境的任务风险 & 代理损失下降不等于目标风险下降 \\
修改局部性 & 按有效请求确定的保留答案 & 无关实体、关系、历史时间及任务 & 参数改动少不等于语义影响小 \\
语义改写泛化 & 与请求含义一致的正确答案 & 等义改写及经核对的跨语言表达 & 原提示成功不等于不同表达均有效 \\
关联一致性 & 已核实依赖及新事实应得结论 & 有依据的依赖推理、多跳关系与版本调用 & 改写成功不等于关联推理正确 \\
能力保留 & 旧任务阶段结果或历史最佳值 & 各阶段仍有效的旧任务表现 & 旧任务不退化不等于后续仍可学习 \\
可塑性 & 同样新证据下的可比起点 & 同预算学习过程及达到阈值的代价 & 约束变弱、参数重置不等于学习恢复 \\
训练影响消除 & 剩余数据的指定学习过程 & 删除对象、剩余任务、参照距离及保证条件 & 拒答、有限输出接近、攻击失败不等于普遍删除保证 \\
\bottomrule
\end{tabularx}
\caption{验收主张、参照和观测对象的对应。各行有独立条件，同一版本可部分通过而其他项失败。保留与可塑性之外，知识积累还需控制相关历史经验有无。表不将各项合并成可相互抵消的总分。}
\label{tab:ket-update-validation}
\end{table*}

\section{知识的持续维护}
\label{sec:ket-update-maintenance}

雨季版本验收后，温室可能重新出现，下一份任命又可能替代当前负责人，记忆也会继续增长。一次通过只对应当时条件，维护需跨轮保存触发证据、预算、有效请求和恢复状态。单轮操作仍调用第~\ref{sec:ket-update-parameters}至第~\ref{sec:ket-update-external}节；以下处理何时调用、如何安排以及前后版本是否相容。

\subsection{更新触发条件的监测}
\label{sec:ket-update-monitoring}

季节更替可能让亮度和物种频率缓慢移动，换相机可能立即改变分辨率、颜色响应或深度单位，临时故障则可能只破坏若干批输入。三者都可触发告警，却须分别调查持续变化、接口变更和短暂异常。先辨认变化来源，再决定持续观察、复核接口或隔离故障，而不直接训练全模型。

输入统计反映观测变化，熵和预测频率反映响应变化，可靠标签才直接揭示任务错误。预测甲比例突增，可能因甲更多，也可能因模型塌缩；亮度变化可能无害，也可能丢纹理。监测应保留这些候选解释，并核对预处理、数据来源和标签。统计越阈值触发调查或可恢复应对，不生成新的正确标签。

标签延迟时，先绑定输入窗口、模型和原预测，必要时暂停高风险更新或增加人工复核，再等反馈决定是否重估规律。后来标签须与预测时的事实时间及计分规则对应，不能用改后的标签口径直接评价旧要求。相关窗口反复越界，也不必是多个独立变化。

触发规则还要在同一事件流上比较误报、漏报与延迟。低阈值可能增加不必要训练，高阈值可能让失效持续；连续越界或多信号核对减少部分误报，也可能延迟突发处理。无任务边界时，窗口及重置条件须预先定义。若已有适用组件，识别上下文后路由即可，不必总生成新参数；具体监测调度归\BookPartRef{part:systems}{卷12《系统》第1部分}。

\subsection{更新节奏与长期预算的安排}
\label{sec:ket-update-budget}

给定同一经验流，可逐次小步更新、攒批更新或按事件启动。小步减少等待，却可能连续放大单一证据的偏差；攒批利用更多观测，增加等待；事件式取决于告警能否找到应更新时点。比较要记录处理次数、标签延迟及停机，而不把“在线”理解为每条只训练一次，回放与重复遍历都应计数。

持续测试时适应CoTTA（continual test-time adaptation）让教师参数取学生的指数滑动平均，以平滑单步变化；源置信度不足时还平均多种增广预测，供学生匹配。同时随机恢复部分源参数以限制偏离。平均可能伴随长期错误，恢复也可能抵消当前适应，均需任务反馈检验。若用于带编辑或撤回的系统，旧源参数可能重新带入失效要求，必须换成合规基准或禁用该恢复；这是本章对其维护适用性的限制，不是原方法已自动处理这些请求\scite{wang2022cotta}。

温室回归时，复测可用能力，并计入恢复样本、更新次数和等待。保存旧模型可能减少重学，增加权重存储；小记忆省模型副本，却需要重新优化，压缩还可能失去稀有覆盖。历史材料持续入库时也须去重、补延迟标签和处理有效期，失效当前负责人答案不能无条件回放，撤回记录不能重新入选。

为约束这种增长，将不重叠的长期存储对象归入五项并求和：
\[
B_{\mathrm{total}}
=B_{\mathrm{model}}+B_{\mathrm{memory}}+B_{\mathrm{state}}
+B_{\mathrm{external}}+B_{\mathrm{recovery}}.
\]
五项依次为模型路由、样本及输出、优化器统计、记录索引和恢复快照，单位须一致，共享对象只计一次。教学上取$(400,250,100,100,150)$，总额1000。新分支需80单位，若其他项不变，只能让记忆从250降至170；不接受覆盖损失时，就须回收、压缩其他对象或拒绝扩容。

按这五项顺序，容量让渡可直接复算为
\[
(400,250,100,100,150)
\ \longrightarrow\
(480,170,100,100,150),
\qquad B_{\mathrm{total}}=1000.
\]
模型增加80、记忆减少80，总额不变。后来新材料仍须在170单位内替换，不能另设不入账的池。

入库时先比较保留价值，再替换旧记录及其标签、特征和概率缓存。类别增多会使logit变宽，照片数固定也不保证字节固定。每类固定内存、总内存固定与模型数固定是不同限制。活跃组件、检索等待、增广次数、更新频率和恢复重算也须累计，不能只报单轮训练耗时。

\subsection{知识版本与有效请求的跟踪}
\label{sec:ket-update-versions}

每次回答应能追到当时适用的依据，而不只追到权重版本。维护记录保存主体、请求来源、区间、修改对象、依赖和验证结果，也保留未决问题。新请求改变条件后，应标记受影响结论待复查；旧验收只对原请求集合和组件组合有效。

续用教学任命，乙从7月1日起，丙从9月1日起。第二次更新将乙区间收束为$[\text{7月1日},\text{9月1日})$，7月15日仍答乙，不随当前任命变丙。另一来源若称9月由丁代理，要核对机构、权限和时间，不能按到达先后覆盖。未决冲突也限制其签批依赖结论，而不只限制直接负责人回答。

请求语义独立，参数操作却可能干扰。共享层再次写入会改变第一次依赖的键或下游表示，因此不同实体也须跨阶段复测。保存批量、长度和依赖复杂度有助区分干扰来源。旧秩一增量虽可存档，网络后来训练后再减它，并不还原全部后续变化；可逆性应由状态依赖与恢复验证判断。

为重建一次响应，可将所用状态标为版本元组：
\[
\begin{aligned}
v=(&v_{\mathrm{data}},v_{\mathrm{prep}},v_{\mathrm{embed}},
v_{\mathrm{train}},v_{\mathrm{model}},v_{\mathrm{eval}},\\
&v_{\mathrm{route}},v_{\mathrm{deploy}},v_{\mathrm{record}},
v_{\mathrm{index}},v_{\mathrm{cache}}),
\end{aligned}
\]
前八项依次绑定数据、预处理、特征生成、训练配置、模型、评价、路由和部署，后三项绑定记录、索引、缓存；另关联有效请求集合$\mathcal E_{\mathrm{active}}$与查询时间。同权重配不同索引可能读到不同证据，同数据经不同预处理也可能形成不同训练和检索结果。图~\ref{fig:ket-update-versions}分别画模型与记录版本，表明恢复不是只换一个文件；这些标识需要指向可追踪内容，不能只有版本字符串。

\begin{figure*}[htbp]
\centering
% Teaching versions: e is an active fact edit; u is a withdrawn record.
% Dotted rollback path is a candidate only, never a direct deployment edge.
\begin{tikzpicture}[
  x=1mm,y=1mm,font=\figurefont\small,
  box/.style={rectangle,draw=BookInk,fill=BookPaper,line width=0.7pt,
    align=center,inner xsep=2mm,inner ysep=2mm,minimum height=13mm},
  flow/.style={draw=BookInk,line width=0.8pt,-{Latex[length=2mm]},
    rounded corners=1mm},
  dependency/.style={flow,dashed,draw=BookBlue},
  candidate/.style={flow,densely dotted,draw=BookAmber,line width=1pt}
]
  \node[anchor=west,text=BookInk] at (6,118) {模型版本：按时间更新};
  \node[box,text width=29mm,fill=FigureModel!10] (m0) at (24,101)
    {$M^{(0)}$\\未满足$e,u$};
  \node[box,text width=29mm,fill=FigureModel!10] (m1) at (82,101)
    {$M^{(1)}$\\已落实$e$};
  \node[box,text width=29mm,fill=FigureOutput!10] (m2) at (140,101)
    {$M^{(2)}$\\已处理$u$};
  \draw[flow] (m0) -- node[above,font=\figurefont\footnotesize] {修正$e$} (m1);
  \draw[flow] (m1) -- node[above,font=\figurefont\footnotesize] {撤回$u$} (m2);

  \node[anchor=west,text=BookInk] at (6,81) {外部事实条目版本：保留有效历史};
  \node[box,text width=29mm] (d0) at (24,65)
    {$D^{(0)}$\\旧事实范围};
  \node[box,text width=29mm] (d1) at (82,65)
    {$D^{(1)}$\\新旧时间区间};
  \node[box,text width=29mm] (d2) at (140,65)
    {$D^{(2)}$\\处理撤回副本};
  \draw[flow] (d0) -- node[above,font=\figurefont\footnotesize] {$e$同步} (d1);
  \draw[flow] (d1) -- node[above,font=\figurefont\footnotesize] {$u$同步} (d2);
  \node[box,text width=48mm,fill=FigureLoss!10] (check) at (55,18)
    {重新核对全部有效请求\\补做$e$及$u$的合规处理};
  \node[box,text width=42mm,fill=FigureOutput!10] (restored) at (136,18)
    {恢复候选$M^{\mathrm R}$\\验收通过后才能启用};
  \draw[candidate] (m0.west) -- (2,101) -- (2,18) -- (check.west);
  \node[anchor=west,text=BookAmber,font=\figurefont\footnotesize]
    at (8,40) {旧快照只作候选来源};
  \draw[flow] (check) -- (restored);
  \draw[dependency] (d2.south) -- (140,42) -| (restored.north);
  \node[anchor=east,text=BookBlue,font=\figurefont\footnotesize]
    at (136,42) {当前事实条目相容};
  \node[align=center,text=BookMuted,font=\figurefont\footnotesize,
    text width=156mm] at (82,-3)
    {实线：版本更新；虚线：相容依赖；点线：待审核恢复路径};
\end{tikzpicture}

\caption{模型与记录的版本依赖及恢复候选，均为教学构造。实线表示更新，虚线表示当前记录相容依赖，点线仅为待审核恢复路径。$e$是有效编辑，$u$是撤回记录；旧模型须补足这些要求，形成候选$M^{\mathrm R}$并验收后才可启用。外部副本处理与参数影响消除仍是两项任务。}
\label{fig:ket-update-versions}
\end{figure*}

\FloatBarrier
\subsection{异常恢复与回退结果的复核}
\label{sec:ket-update-rollback}

当前适应失败时，旧快照提供恢复起点，却未必满足今天的要求。先确定要恢复的能力与可用状态，再核对有效编辑、行为约束和撤回集合；快照缺少要求时，补做操作或从合规材料重建。权重、路由和外部链路须相容，无法验证时不能直接启用。回退本身是一项待验收变更。

教学上设$M^{(0)}$温室效果可用，却带后来撤回记录$u$的影响；$M^{(1)}$完成任命编辑，$M^{(2)}$又处理$u$。后续适应退化后回到$M^{(0)}$，可能同时恢复温室能力、旧负责人错误和撤回影响。须由合规快照出发，或对旧快照补做编辑与删除，形成$M^{\mathrm R}$再验收。若删除所需状态已失，低成本补做并不一定可行。

外部库恢复也须保留当前责任。旧任命可服务历史，却不能覆盖现任；撤回缓存与索引不能因备份恢复重新开放。模型和库的有效请求不一致，会制造新错误。恢复测试因此同时覆盖目标能力、历史时间、无关响应、撤回参照与版本组合，通过后继续监测回归域和后续学习。

代理持续改善而真实反馈恶化、来源冲突未决、缺合规恢复路径，或预算不能维持必需能力时，可以暂停相关更新并限制服务。暂停不等于修复，应记录可提供能力、待决问题与重新评估条件。日志、存储和部署实现归系统卷，本章规定当前行为必须追到明确版本、条件和证据。

\FloatBarrier
\section{本章小结}
\label{sec:ket-update-summary}

面对雨季错误、负责人任命或记录撤回，先确定哪种关系需要变，再准备能够辨认它的证据。频率校正依赖保留的条件分布，观测更换依赖可靠对应与可用信息，事实修正依赖身份、时间和推理规则，机器遗忘依赖规定的剩余数据学习参照。请求可以组合，证据却不能相互替代。

参数变化由经验贡献、目标约束、可变范围和计算共同决定，结构配置及外部记录提供另外两类修改对象。回放覆盖输入，行为约束保护输出，二次罚项限制参数，它们可以组合但不等价；冻结旧组件仍须保护正确路由，改好记录仍须同步调用。保留集合也要排除合法修正和撤回已经否定的旧要求，允许有依据的依赖传播。

是否接受更新，需要分别检查当前风险、修改范围、旧能力、后续学习和删除保证，不能靠一个平均分。持续维护再把这些要求与有限预算、有效请求及版本依赖一起保存，使恢复候选也受相同约束；旧域分数更好不能授权重新引入已撤回影响。用途扩展由第~\ref{chap:ket-transfer-adaptation}章确定新任务，所得能力进入服务后仍须遵守这里的保留、验收与维护条件。
